\documentclass[11pt]{article}

\usepackage[T1]{fontenc}
\usepackage[utf8]{inputenc}
\usepackage{lmodern}
\usepackage{amsmath,amssymb,amsthm,mathtools}
\usepackage[a4paper,margin=28mm]{geometry}
\usepackage{microtype}
\usepackage{xurl}
\usepackage[hidelinks]{hyperref}
\usepackage{enumitem}
\usepackage{booktabs}
\usepackage{array}

\newtheorem{theorem}{Theorem}[section]
\newtheorem{lemma}[theorem]{Lemma}
\newtheorem{proposition}[theorem]{Proposition}
\newtheorem{corollary}[theorem]{Corollary}
\newtheorem{definition}[theorem]{Definition}
\newtheorem{remark}[theorem]{Remark}

\newcommand{\N}{\mathbb N}
\newcommand{\WIN}{\ensuremath{\mathrm{WIN}}}
\newcommand{\LOSS}{\ensuremath{\mathrm{LOSS}}}
\newcommand{\DRAW}{\ensuremath{\mathrm{DRAW}}}
\newcommand{\odd}{\operatorname{odd}}
\newcommand{\src}{\operatorname{src}}
\newcommand{\ct}{\operatorname{ct}}
\newcommand{\asl}{\operatorname{asl}}
\newcommand{\Moves}{\operatorname{moves}}

\title{The Two-Player $3n\mathbin{\pm}1$ Game:\\
A Binary Normal Form and a Well-Founded No-Draw Proof}
\author{Grisha Pochuev\\
  \small\href{mailto:n_854@mail.ru}{n\_854@mail.ru}\\
  \small\url{https://github.com/Grisha-Pochuev/3n-plus-minus-1-game}}
\date{Audit-hardened synthesis version 6.3, 11 August 2026}

\begin{document}
\maketitle

\begin{abstract}
We study the two-player $3n\pm1$ game on positive odd integers. Earlier
versions developed the binary normal form, constant-tail coordinates, source
embedding, factor-removal diamonds and proof-height tokens, but hostile audits
found two genuine global defects: a coefficient decrease was once confused
with a source decrease, and a retained source from before a multi-$A$ streak
was compared with a later $B$ return for which that provenance no longer held.
Both shortcuts remain permanently withdrawn.

This version closes the two routing obligations isolated by the audit.  The
new bridge is an exact raw-selector lemma: after a canonical source transition,
the next phase selects the precise ordinary child already exposed by the
preceding DRAW diamond.  Hence a valuation-two recycle can seed its first
finite WIN token only once and every later free recycle is tied to a strict
proof-tree descendant; no pre-streak numerical anchor is used.  The
factor/high-return side is closed separately by retaining proof tokens through
the full semantic case split.  Short high returns spend a finite descendant,
valuation at least seven decreases a two-token multiset, every marked tail of
length at least three replaces its LOSS token by an actual lower WIN child,
and the two short tail rows remain exact lifts over an already retained token.
A lexicographic multiset rank then rules out an infinite marked DRAW
normalization.  The machine checks remain supporting regressions rather than an
end-to-end formal proof.
\end{abstract}

\tableofcontents

\section{The game and the theorem}

The current state is a positive odd integer $n$. For $n>1$ the player to move
chooses one of
\[
  3n-1,\qquad 3n+1,
\]
and divides the chosen value by the largest possible power of $2$, leaving an
odd integer. If the resulting odd integer is $1$, the player who made the move
wins immediately. Players alternate with perfect information.

A position is called \WIN\ for the player to move if at least one move leads to
\LOSS. It is called \LOSS\ if every legal move leads to \WIN. These two classes
are the inductively generated finite-remoteness classes. Any position not in
either class is called \DRAW.

\begin{theorem}[No-DRAW theorem]\label{thm:main}
Every positive odd starting position has finite remoteness. Equivalently, the
$3n\pm1$ game has no \DRAW\ positions and optimal play reaches $1$ after
finitely many moves.
\end{theorem}

The theorem is not a statement about arbitrary play. There are legal cycles,
such as
\[
  5\longrightarrow 7\longrightarrow 5,
\]
under suitable choices by the players.

\section{Binary conjugation}

Write an odd state as $n=2m+1$ and define
\[
  F(m)=\left\lceil\frac{3m}{2}\right\rceil .
\]
For a nonnegative integer $x$, let $R(x)$ be the integer obtained by deleting
the maximal alternating suffix of the binary expansion of $x$. Thus
$R(1001_2)=10_2$ and $R(10101_2)=0$.

\begin{proposition}[First normal form]\label{prop:first-normal}
For $m>0$, the two children in $m$-coordinates are exactly
\[
  F(m),\qquad F(R(m)).
\]
\end{proposition}

\begin{proof}
One has
\[
  3(2m+1)-1=2(3m+1),\qquad
  3(2m+1)+1=2(3m+2).
\]
Exactly one of $3m+1$ and $3m+2$ is odd. In the other expression, each
additional factor of two removed corresponds to crossing one additional
alternating bit at the end of the binary expansion of $m$. The process stops
at the first repeated adjacent bit, which is precisely the deletion defining
$R(m)$. Returning to $m$-coordinates gives the displayed pair.
\end{proof}

Since both children lie in the image of $F$, we represent the state $F(q)$ by
$q$. The conjugated game has terminal state $0$ and moves
\[
  A(q)=\left\lceil\frac{3q}{2}\right\rceil,
  \qquad
  B(q)=R(A(q)).
\]

\begin{proposition}[Strict contraction]\label{prop:B-contracts}
For every $q>0$,
\[
  B(q)<q.
\]
\end{proposition}

\begin{proof}
Deleting a nonempty binary suffix gives
\[
  B(q)\le \left\lfloor\frac{A(q)}2\right\rfloor
       \le \left\lfloor\frac{3q+1}{4}\right\rfloor<q
\]
for $q\ge2$, while $B(1)=0$.
\end{proof}

Thus all unbounded difficulty comes from the expanding branch $A$ and from the
unbounded alternating suffix removed by $B$.

\section{Constant-tail coordinates and coefficient sources}

For a positive odd integer $a$, an exponent $r\ge1$, and a phase
$e\in\{0,1\}$, put
\[
  Q_r^e(a)=a2^r-e.
\]
This is a binary word whose final constant block has length $r$ and consists
of the bit $e$.

\begin{lemma}[Long-tail recurrence]\label{lem:long-tail}
For every positive odd $a$, phase $e$, and $r\ge3$,
\[
  A(Q_r^e(a))=Q_{r-1}^e(3a),\qquad
  B(Q_r^e(a))=Q_{r-2}^e(3a).
\]
Moreover
\[
  B(Q_1^e(a))=B(Q_2^e(a)).
\]
\end{lemma}

\begin{proof}
For $r\ge3$, the expanding child still ends in at least two equal bits; hence
its maximal alternating suffix consists only of the final bit. This gives both
long-tail identities. At the boundary, the binary words of $3a-e$ and $6a-e$
differ by appending one bit that extends the same alternating suffix, so the
same prefix remains after deletion.
\end{proof}

Define the embedded three-free coefficient
\[
  J(s)=2A(s)+1.
\]
Every positive odd coefficient $a$ has a unique factorization
\[
  a=3^kJ(s),\qquad k\ge0.
\]
Indeed, remove all factors of $3$ and call the remaining odd integer $c$.
Then $c\not\equiv0\pmod3$ and $(c-1)/2$ is $0$ or $2$ modulo $3$.
The image of
\[
 A(2t)=3t,\qquad A(2t+1)=3t+2
\]
is exactly the set of nonnegative integers not congruent to $1$ modulo $3$,
and $A$ is injective. Hence there is a unique $s$ with
$A(s)=(c-1)/2$, equivalently $c=J(s)$. Unique factorization supplies the
unique exponent $k$. The integer $s$ is called the coefficient source of
$a$ and is denoted $\src(a)$.

\begin{lemma}[Source boundary]\label{lem:source-boundary}
For every $s>0$ there is a phase
\[
  \alpha(s)=1-\left(\left\lfloor\frac{s}{2}\right\rfloor\bmod2\right)
\]
such that
\[
  \odd(3J(s)+1-2\alpha(s))=J(A(s))
\]
with $2$-adic valuation one, while the opposite phase selects $J(B(s))$ with
valuation at least two.
\end{lemma}

\begin{proof}
Put $u=A(s)$. Then $J(s)=2u+1$ and
\[
  3J(s)+1-2e=2(3u+2-e).
\]
The bracket is odd exactly when $e=1-(u\bmod2)$. For $s=2t$ or $2t+1$ one
has $u\bmod2=t\bmod2$, so this phase is exactly $e=\alpha(s)$. If $u$ is
even, then $e=1$ and $3u+1=J(u)$; if $u$ is odd, then $e=0$ and
$3u+2=J(u)$. Thus the valuation-one child is $J(A(s))$. The opposite sign is
the other original $3n\pm1$ child of the odd state $J(s)$; under the
conjugacy of Proposition~\ref{prop:first-normal} that child is $J(B(s))$.
Since it is not the valuation-one branch, its valuation is at least two.
\end{proof}

Multiplying a constant-tail coefficient by $3$ preserves its coefficient
source. Passing through a factor-free signed boundary with coefficient $J(s)$
therefore exposes exactly one ordinary $A/B$ move on the source.

For a positive state $q$, write $\kappa(q)$ for the odd coefficient in its
canonical constant-tail coordinates.

\begin{lemma}[Coefficient pullback]\label{lem:coeff-pullback}
For every positive odd $w$,
\[
  \kappa(3w)=J(R(w)),\qquad
  \kappa(3w-1)=J(R(2w-1)).
\]
\end{lemma}

\begin{proof}
Let $m\ge1$ be the length of the maximal alternating suffix of the odd word
$w$, and put $z=R(w)$. If $z>0$, maximality says that the leading bit of the
suffix equals the last bit of $z$. Because the suffix ends in $1$, its value is
\[
 t_m=\begin{cases}
 (2^m-1)/3,&m\text{ even},\\
 (2^{m+1}-1)/3,&m\text{ odd}.
 \end{cases}
\]
Thus $w=2^mz+t_m$. If $m$ is even, $z$ is even and
\[
  3w+1=2^m(3z+1)=2^mJ(z).
\]
If $m$ is odd, $z$ is odd and
\[
  3w+1=2^m(3z+2)=2^mJ(z).
\]
If $z=0$, the whole word is alternating; then $m$ is odd and
$3w+1=2^{m+1}$, whose odd part is $J(0)=1$. Hence in all cases
$\kappa(3w)=\odd(3w+1)=J(R(w))$.

For the second identity set $v=2w-1$, which is odd. Since
$3v+1=2(3w-1)$, the first identity applied to $v$ gives
\[
 \kappa(3w-1)=\odd(3w-1)=\odd(3v+1)=J(R(v))=J(R(2w-1)).
\]
\end{proof}


\section{Finite outcomes, witnesses, and ranks}

The terminal state $0$ is \LOSS\ with height $h(0)=0$. A finite \WIN\ state has
height
\[
  h(x)=1+\min\{h(y):y\text{ is a \LOSS\ child of }x\},
\]
and a finite \LOSS\ state has height
\[
  h(x)=1+\max\{h(y):y\text{ is a child of }x\}.
\]
A state whose outcome has been proved but which is being carried only to
justify a later incidence is called a \emph{routing witness}. A
\emph{retained proof token} is a finite state for which the proof has also
established the comparison needed by the rank: it either seeds an empty inner
token set at a one-shot seed transition, or replaces a previously retained
token by certified descendants of strictly smaller height. This distinction is
important: a newly discovered finite state is never silently appended to a
ranked multiset.

\begin{lemma}[Multiset token rank]\label{lem:token-rank}
For a finite multiset $M$ of proof heights define
\[
  \Phi(M)=\mathop{\#}_{h\in M}\omega^h,
\]
the natural ordinal sum of the monomials $\omega^h$. Replacing one token of
height $H$ by any finite multiset of certified descendants of heights strictly
below $H$ strictly decreases $\Phi$.
\end{lemma}

\begin{proof}
In Cantor normal form the coefficient of $\omega^H$ decreases by one, while
all inserted terms have lower exponent. No finite collection of lower terms
can compensate for the loss at exponent $H$.
\end{proof}

Every positive state $q$ has unique canonical constant-tail coordinates
$q=Q_r^e(a)$ with $a$ odd and $r\ge1$. If
$a=3^kJ(s)$, define the \emph{state source} $\rho(q)=s$.

\begin{lemma}[Universal state-source bound]\label{lem:source-bound}
For every $q>0$,
\[
  \rho(q)\le \frac{q-1}{6}.
\]
\end{lemma}

\begin{proof}
Write $q=Q_r^e(3^kJ(s))$. Since $r\ge1$ and $e\le1$,
$q\ge2J(s)-1$. Moreover
$J(s)=2\lceil3s/2\rceil+1\ge3s+1$. Hence
$q\ge6s+1$.
\end{proof}

\begin{lemma}[Boundary coefficient descent]\label{lem:boundary-coeff}
Let $a$ be positive and odd and
\[
 p=B(Q_1^e(a))=B(Q_2^e(a)).
\]
If $p>0$, then its canonical coefficient satisfies
\[
  \kappa(p)\le\frac{3a+1}{4}.
\]
In particular $\kappa(p)<a$ for $a\ge3$.
\end{lemma}

\begin{proof}
The common child is $p=R(3a-e)$. Deleting a nonempty suffix gives
$p\le\lfloor(3a-e)/2\rfloor$. For every positive integer $y$, the odd
coefficient in its canonical constant-tail coordinates is at most $(y+1)/2$.
Therefore
\[
 \kappa(p)\le\frac{p+1}{2}\le\frac{3a+1}{4}.
\]
The final inequality is strict relative to $a$ for $a\ge3$.
\end{proof}

\section{Side identities and typed boundary objects}

The finite routing system uses two elementary binary identities. We record
them here because they are the only places where an exceptional residue class
enters.

\begin{lemma}[Ordinary side relation]\label{lem:side}
Put $S(q)=B(q)$ and $T(q)=B(A(q))$. If
\[
  q\not\equiv1,3,12,14\pmod{16},
\]
then
\[
  T(q)\in\{A(S(q)),B(S(q))\}.
\]
If $q$ is in one of the four exceptional classes, then $A(q)$ is not
exceptional.
\end{lemma}

\begin{proof}
Let $x=A(q)$ and let $k$ be the length of the maximal alternating suffix of
$x$. Direct reduction of $A(q)$ modulo $8$ shows that $k\ge3$ is possible
exactly for $q\equiv1,3,12,14\pmod{16}$. Outside these classes $k=1$ or
$2$. Writing $x$ as $S(q)$ followed by that one- or two-bit alternating word,
the formulas
\[
  A(2u)=3u,\qquad A(2u+1)=3u+2
\]
show that $A(x)$ is obtained from $A(S(q))$ by appending the corresponding
one- or two-bit alternating word. Deleting its maximal alternating suffix
therefore leaves either $A(S(q))$ or $B(S(q))$. The four exceptional residue
classes map under $A$ to $2,5,10,13\pmod{16}$, none of which is exceptional.
\end{proof}

\begin{corollary}[No four consecutive wins]\label{cor:no-four-win}
For no $q$ can all four states
\[
  q,A(q),A^2(q),A^3(q)
\]
be \WIN.
\end{corollary}

\begin{proof}
If $q$ and $A(q)$ are both \WIN, then $B(q)$ must be \LOSS; the same holds at
the next two positions. If $q$ is ordinary, Lemma~\ref{lem:side} makes
$B(A(q))$ a child of the \LOSS\ state $B(q)$, impossible. Hence $q$ must be
exceptional. Applying the same argument to $A(q)$ forces $A(q)$ exceptional,
contrary to Lemma~\ref{lem:side}.
\end{proof}

\begin{lemma}[Alternating lift identity]\label{lem:alt-lift}
For every positive integer $z$ there are $m\ge1$ and
$\delta=1-(z\bmod2)$ such that
\[
  3z+1=Q_m^\delta(J(R(z))).
\]
If $R(z)>0$, one may take $m$ equal to the length of the maximal alternating
suffix of $z$; if $R(z)=0$, one takes one more than the binary length of $z$.
\end{lemma}

\begin{proof}
First suppose $r=R(z)>0$, let $k$ be the alternating-suffix length and put
$\epsilon=r\bmod2$. Then
\[
 z=2^k r+t_{k,\epsilon},\qquad
 t_{k,0}=\Big\lfloor\frac{2^k}{3}\Big\rfloor,
 \quad t_{k,1}=\Big\lfloor\frac{2^{k+1}}{3}\Big\rfloor.
\]
Since $J(r)=3r+1+\epsilon$, direct substitution gives
\[
 3t_{k,\epsilon}+1+\delta=2^k(1+\epsilon),
\]
the two parity cases being exactly the two remainders of a power of $2$
modulo $3$. Hence
\[
 3z+1+\delta=2^kJ(r),
\]
which is the required identity with $m=k$. If $R(z)=0$, the whole binary word
of $z$ is alternating and starts with $1$, so
$z=t_{k,1}$. The same calculation gives
$3z+1+\delta=2^{k+1}=2^{k+1}J(0)$.
\end{proof}

\begin{lemma}[Exceptional child-pair identity]\label{lem:exceptional-side}
Let $q\equiv1,3,12,14\pmod{16}$ and put $\ell=B(q)$.  Then for some
$m\ge1$ and $\delta\in\{0,1\}$ the two children of $A(q)$ are exactly
\[
  Q_m^\delta(J(\ell)),\qquad Q_{m+1}^\delta(J(\ell)).
\]
Consequently, if in addition $q$ is \WIN, $A(q)$ is \DRAW, and
$\ell$ is \LOSS, this adjacent pair contains a continuing \DRAW.
\end{lemma}

\begin{proof}
Write $q=16t+c$ with $c\in\{1,3,12,14\}$. The four direct exceptional
formulas are
\[
\begin{array}{c|c|c|c}
c&\ell&B(A(q))&A^2(q)\\ \hline
1&R(6t)&18t+1&36t+3\\
3&R(6t+1)&18t+4&36t+8\\
12&R(6t+4)&18t+13&36t+27\\
14&R(6t+5)&18t+16&36t+32.
\end{array}
\]
In every row put respectively $z=6t,6t+1,6t+4,6t+5$. Then
$\ell=R(z)$ and $B(A(q))=3z+1$. By Lemma~\ref{lem:alt-lift},
\[
  B(A(q))=Q_m^\delta(J(\ell))
\]
for some $m\ge1$. The last column is $2B(A(q))+\delta$, so it equals
$Q_{m+1}^\delta(J(\ell))$. These identities use no outcome hypothesis and
are precisely the two children of $A(q)$.  If $A(q)$ is \DRAW, neither
child can be \LOSS and at least one is \DRAW, proving the consequence.
\end{proof}

For a source $x$ and phase $e$, let $O(x,e)$ denote the following exact
\emph{DRAW obligation}:
\begin{enumerate}[label=(\alph*)]
\item $Q_1^e(J(x))$ is \DRAW; or
\item there is a \DRAW\ state whose two children are exactly
      $Q_2^e(J(x))$ and $Q_1^e(J(x))$.
\end{enumerate}
This definition introduces no outcome assumption beyond the displayed
\DRAW\ state.

\begin{lemma}[Hidden parent for a typed three/two frame]\label{lem:hidden-parent}
Let $a$ be positive, odd and not divisible by $3$, and suppose
\[
  a\equiv1+e\pmod3.
\]
Put $H=Q_3^e(a)$, $G=Q_2^e(a)$ and
\[
  K=\frac{2H+e-1}{3}.
\]
Then $K$ is an integer and
\[
  \Moves(K)=\{H,G\}.
\]
\end{lemma}

\begin{proof}
For $e=0$ the congruence gives $H\equiv2\pmod3$ and
$K=(2H-1)/3$; for $e=1$ it gives $H\equiv0\pmod3$ and $K=2H/3$.
These are exactly the two inverse branches of the injective map
$A(q)=\lceil3q/2\rceil$. The last three bits of $H$ are equal, so the
contracting child removes only the last bit and equals $G$.
\end{proof}

\section{Fixed-tail factor removal and entry typing}

The next two lemmas isolate the only raw entry phenomenon that is not already
an obligation or a factor frame.

\begin{lemma}[Fixed-tail factor-removal diamond]\label{lem:fixed-tail-factor}
Let $c$ be positive and odd, $e\in\{0,1\}$, and $r\ge2$. Put
\[
  X=Q_r^e(3c),\qquad G=Q_r^e(c),\qquad H=Q_{r+1}^e(c),\qquad Y=A(G).
\]
Then
\[
  \Moves(H)=\{X,Y\}.
\]
If $X$ is \DRAW, either $G$ is \DRAW, or the diamond determines a finite
routing witness together with its exact incidence to $G,H,X,Y$ and the other
child of $G$.
\end{lemma}

\begin{proof}
Since $r+1\ge3$, Lemma~\ref{lem:long-tail} gives $A(H)=X$. For $r\ge3$ it
also gives $B(H)=Q_{r-1}^e(3c)=A(G)$; for $r=2$ the boundary identity gives
the same equality. Thus $B(H)=Y$.

Because $X$ is \DRAW, $H$ is not \LOSS. If $H$ is \WIN, then its other
child $Y$ is necessarily \LOSS, and this is the required finite witness.
Assume instead that $H$ is \DRAW. Then $Y$, being a child of a DRAW, is not
LOSS. If $G$ is \LOSS, then both children of $G$, including $Y=A(G)$, are
WIN. If $G$ is \WIN, at least one child of $G$ is LOSS; since $Y$ is not LOSS,
that witness is the other child $B(G)$. These are all outcome orientations.
In every finite branch the witness and the parent-child incidence just
displayed are known exactly; no token-rank comparison is asserted at this
stage.
\end{proof}

\begin{definition}[Entry controls]\label{def:entry-controls}
A \emph{boundary entry} is an obligation $O(x,e)$ or, more generally, a
factor-free adjacent frame
\[
 \{Q_{m+1}^e(J(x)),Q_m^e(J(x))\},\qquad m\ge1,
\]
known to contain a continuing \DRAW. Its remaining tail length is recorded as
$\tau$. The typed three/two frame satisfying Lemma~\ref{lem:hidden-parent} is
a distinguished $m=2$ boundary constructor.
A \emph{loss-sibling entry} is one of the finitely many local constructors
in which a continuing \DRAW\ is accompanied by a finite \WIN/\LOSS\ routing
witness with a proved parent/child incidence: the fixed-tail diamond of Lemma~\ref{lem:fixed-tail-factor}, the exponent-one predecessor diamond of Lemma~\ref{lem:exp-one-lift}, an ordinary side diamond of Lemma~\ref{lem:side}, the exceptional side attachment of Lemma~\ref{lem:exceptional-side}, or the hidden-parent frame of Lemma~\ref{lem:hidden-parent}. The constructor records its remaining
constant-tail/factor exponent as the local counter $\tau$.
\end{definition}

\begin{lemma}[Entry typing]\label{lem:entry-typing}
Every finite branch produced by Lemma~\ref{lem:fixed-tail-factor} is a
loss-sibling entry in the sense of Definition~\ref{def:entry-controls}; no
ranked multiset is changed in order to make this classification.
\end{lemma}

\begin{proof}
The proof of Lemma~\ref{lem:fixed-tail-factor} gives the finite witness and
its exact incidence to $G,H,X,Y$ and the other child of $G$. These are
precisely the data of the first loss-sibling constructor. Classification is
therefore immediate; subsequent tail/factor normalization belongs to the
typed routing relation and is ranked there by $\tau$ and, when available,
by descendant-token comparisons.
\end{proof}

\begin{lemma}[Exponent-one lift]\label{lem:exp-one-lift}
Let $a$ be positive and odd and $k\ge1$. If
$R_k=Q_1^e(3^ka)$ is \DRAW, then
\[
  P_{k-1}=Q_2^e(3^{k-1}a)
\]
is either \DRAW\ or \WIN; in the latter case its other child is a finite
\LOSS\ witness.
\end{lemma}

\begin{proof}
The boundary recurrence gives $A(P_{k-1})=R_k$. A state with a \DRAW\ child
cannot be \LOSS. If $P_{k-1}$ is \WIN, the \DRAW\ child cannot witness the
win, so its other child is \LOSS.
\end{proof}

\begin{proposition}[Arbitrary factorful exponent-one entry]\label{prop:factorful-entry}
Let $a=J(s)$ and $k>0$. If $Q_1^e(3^ka)$ is \DRAW, then after finitely many
outcome-compatible local steps one obtains either a typed entry of
Lemma~\ref{lem:entry-typing}, or the factor-free exponent-two state
\[
  Q_2^e(a)\text{ \DRAW}.
\]
\end{proposition}

\begin{proof}
Apply Lemma~\ref{lem:exp-one-lift}. Its finite branch is a loss-sibling entry by Definition~\ref{def:entry-controls}. Otherwise $Q_2^e(3^{k-1}a)$ is \DRAW. Apply
Lemma~\ref{lem:fixed-tail-factor} with $r=2$ repeatedly. Whenever its finite
branch occurs, Lemma~\ref{lem:entry-typing} supplies a typed entry. If it never
occurs, one factor of $3$ is removed at the same tail exponent at each step,
so after $k-1$ steps the displayed factor-free exponent-two state is reached.
\end{proof}

\begin{proposition}[Complete exponent-one/two boundary reduction]\label{prop:boundary-reduction}
Let $s\ge0$, $k\ge0$, $e\in\{0,1\}$, and let
\[
  Q_r^e(3^kJ(s)),\qquad r\in\{1,2\},
\]
be \DRAW. For $s>0$, after finitely many local steps at least one of the following occurs: a typed entry is obtained, the retained source strictly
decreases, or the factor-free exponent-two state $Q_2^e(J(s))$ is \DRAW.
For $s=0$ the corresponding statement is Lemma~\ref{lem:zero-source}.
\end{proposition}

\begin{proof}
Assume $s>0$. There are four cases. If $r=1,k=0$, the state is the first
alternative of the obligation $O(s,e)$ and is already a boundary entry. If
$r=1,k>0$, use Proposition~\ref{prop:factorful-entry}. If $r=2,k=0$, we are
at the displayed factor-free exponent-two state. Finally, if $r=2,k>0$, apply
Lemma~\ref{lem:fixed-tail-factor} repeatedly with $r=2$: a finite branch is a
loss-sibling entry by Definition~\ref{def:entry-controls}, while otherwise
one factor of $3$ is removed at the same exponent at each step and after $k$
steps the factor-free exponent-two state is reached. These cases are disjoint
and exhaustive.
\end{proof}

\section{Factor-free exponent-two base entry}

Let $x>0$, $a=J(x)$ and
\[
  P=Q_1^e(a),\quad U=Q_2^e(a),\quad
  b=B(P)=B(U),\quad F=A(U)=Q_1^e(3a),\quad g=1-e.
\]
Assume $U$ is \DRAW\ and $x$ is the least coefficient source of any \DRAW.  This is a \emph{factor-free} coefficient $J(x)$.  The common child $b$ cannot be $0$, because $0$ is \LOSS\ and a \DRAW\ has no \LOSS\ child.  By Lemma~\ref{lem:boundary-coeff}, the canonical coefficient of $b$ is strictly below $J(x)$.  After deleting any factor of $3$ its three-free coefficient is smaller still, hence it equals $J(t)$ with $t<x$ because $J$ is strictly increasing.  Therefore the state source of $b$ is below $x$, and $b$ is not \DRAW.  No such inference is made for a coefficient $3^kJ(x)$ with $k>0$.
As a child of a \DRAW\ it is not \LOSS, so
\[
  b\text{ is \WIN},\qquad F\text{ is \DRAW}.
\]
\begin{lemma}[First-factor anchor identity]\label{lem:first-factor}
For the states above,
\begin{equation}\label{eq:first-factor}
  9J(x)+1-2e=2^vJ(b),\qquad v\ge1.
\end{equation}
\end{lemma}

\begin{proof}
Put $a=J(x)$. If $e=0$, then $P=2a$, $A(P)=3a$ and $b=R(3a)$.
Apply the first identity of Lemma~\ref{lem:coeff-pullback} to the odd integer
$w=3a$:
\[
 \odd(9a+1)=\kappa(9a)=J(R(3a))=J(b).
\]
If $e=1$, then $P=2a-1$, $A(P)=3a-1$ and $b=R(3a-1)$. The second identity
of Lemma~\ref{lem:coeff-pullback}, again with $w=3a$, gives
\[
 \odd(9a-1)=J(R(6a-1)).
\]
By the boundary identity of Lemma~\ref{lem:long-tail},
$R(6a-1)=R(3a-1)=b$. Hence in both phases the odd part of
$9a+1-2e$ is exactly $J(b)$, which proves \eqref{eq:first-factor}.
\end{proof}

\begin{proposition}[Base-entry attachment]\label{prop:base-entry}
Every factor-free exponent-two \DRAW\ at a globally minimal source has a
finite continuation which either strictly lowers the retained source or
produces a boundary/loss-sibling typed entry.
\end{proposition}

\begin{proof}
Suppose first that $e=1-\alpha(x)$. The first source-boundary numerator is
divisible by $4$, while
\[
9J(x)+1-2e=3(3J(x)+1-2e)-2(1-2e)\equiv2\pmod4.
\]
Thus $v=1$. Writing $u=A(x)$ gives $b=3u+1$. The two children of the \DRAW\
state $F$ are
\[
  T=Q_1^g(J(b)),\qquad C=R(T),
\]
and direct appended-bit deletion gives
$C\in\{A(b),B(b)\}$. If $T$ is \DRAW, this is the boundary entry
$O(b,g)$. If $C$ is \DRAW, the other child $\ell$ of the \WIN\ state $b$ is a
finite \LOSS\ witness. If $b$ is ordinary, this is the ordinary-side
loss-sibling constructor. If $b$ is exceptional and $C=A(b)$,
Lemma~\ref{lem:exceptional-side} gives an adjacent \DRAW\ frame over
$\ell$, hence a boundary/loss-sibling entry. It remains only the exceptional
orientation $C=B(b)$. In that orientation
$b\le(9x+5)/2$ and the alternating suffix of $A(b)$ has length at least three,
so
\[
 C=B(b)\le\frac{A(b)}8\le\frac{3b+1}{16}.
\]
Lemma~\ref{lem:source-bound} then yields
\[
 \rho(C)\le\frac{C-1}{6}<\frac{3b+1}{96}
 \le\frac{27x+17}{192}<x,
\]
a strict retained-source exit.

Now suppose $e=\alpha(x)$. Lemma~\ref{lem:source-boundary} gives
$3J(x)+1-2e=2J(A(x))$, hence $v\ge2$. If $v\ge4$, then
\[
16J(b)\le9J(x)+1\le27x+19,
\]
and $J(b)\ge3b+1$ implies
\[
  b\le\frac{9x+1}{16}<x.
\]
Thus the continuing \DRAW\ has smaller retained source. If $v=2$, the two
children of $F$ are exactly $Q_2^g(J(b))$ and $Q_1^g(J(b))$, which is the
second form of $O(b,g)$. If $v=3$, they are the typed three/two frame
$Q_3^g(J(b)),Q_2^g(J(b))$. Reducing
\eqref{eq:first-factor} modulo $3$ gives
\[
  J(b)\equiv1+g\pmod3,
\]
so Lemma~\ref{lem:hidden-parent} applies. The hidden parent is a boundary
entry; its finite side, if one is exposed, is a loss-sibling entry. Hence no
untyped base row remains.
\end{proof}

\section{The source-zero branch}

For $k\ge0$, $r\ge1$ and $e\in\{0,1\}$ put
\[
  S_{k,r}^e=Q_r^e(3^k).
\]
These are precisely the constant-tail states with coefficient source zero.

\begin{lemma}[Zero-source boundary normalization]\label{lem:zero-source}
If a source-zero \DRAW\ exists, then after finitely many outcome-compatible
moves one reaches a typed boundary or loss-sibling entry.  No ranked token is
installed and no control reset is used before this finite normalization is
complete.
\end{lemma}

\begin{proof}
For $r\ge3$, Lemma~\ref{lem:long-tail} gives
\[
  \Moves(S_{k,r}^e)=\{S_{k+1,r-1}^e,S_{k+1,r-2}^e\}.
\]
Following a \DRAW\ child strictly lowers $r$, so a \DRAW\ with $r\in\{1,2\}$
is reached after finitely many moves. If $k=0$ already at that boundary, the
four possibilities are explicit:
\[
 Q_1^1(1)=1\ \WIN,\quad Q_2^1(1)=3\ \WIN,\quad
 Q_1^0(1)=2\ \LOSS,\quad Q_2^0(1)=4\ \LOSS.
\]
Hence a boundary \DRAW\ must have accumulated at least one factor of $3$.

If the reached boundary has tail exponent one, put $n=k+1$.  If it has
tail exponent two, its children are the common raw child and
$S_{k+1,1}^e$.  If the raw child is \DRAW, again put $n=k+1$; otherwise the
raw child is \WIN\ and forces $S_{k+1,1}^e$ to be \DRAW, in which case put
$n=k+2$.  Thus in all cases one obtains a boundary \DRAW\ configuration whose
raw/signed exits are governed, for some $n\ge1$, by
\[
  3^n+1-2e=2^jJ(y).
\]
The exact $2$-adic valuation is
\[
\begin{array}{c|cc}
 &n\text{ odd}&n\text{ even}\\ \hline
 e=0&j=2&j=1\\
 e=1&j=1&j=2+v_2(n).
\end{array}
\]
For $j\ge2$ the raw and signed exits form the adjacent factor-free frame
\[
  Q_{j-1}^{1-e}(J(y)),\qquad Q_j^{1-e}(J(y)),
\]
containing a continuing \DRAW; by Definition~\ref{def:entry-controls} this is a boundary entry, with its finite tail length recorded in $\tau$. For $j=1$ one has
\[
  y=\frac{3^{n-1}-1}{2}.
\]
When $e=0$ the raw exit is $A(y)$; when $e=1$ it is $B(y)$. The raw and
signed states are the two exits of the same boundary \DRAW\ configuration.
If the signed exit is \DRAW, it is already the exponent-one form of an
obligation. If the raw exit is the continuing \DRAW\ and the signed exit is
finite, that finite state is necessarily \WIN; choosing a canonical \LOSS\
child of it supplies an exact finite routing witness, hence a loss-sibling
entry. If both exits are \DRAW, use the signed obligation. Thus every row
reaches a typed entry after a finite pre-entry computation; no impossible
``source below zero'' comparison is used.
\end{proof}

\begin{lemma}[$A$-selecting side diamond]\label{lem:a-side}
Let $x>0$, $e=\alpha(x)$, $w=Q_1^e(J(x))$, and $y=A(x)$. Then
\[
  B(w)\in\{A(y),B(y)\}.
\]
\end{lemma}

\begin{proof}
The source-boundary identity gives
$A(w)=Q_1^{1-e}(J(y))=4A(y)+1+e$. For $e=0$ the appended bits are $01$;
for $e=1$ they are $10$. If the first appended bit repeats the last bit of
$A(y)$, deleting the alternating suffix leaves $A(y)$; otherwise the deletion
continues through the maximal alternating suffix of $A(y)$ and leaves
$B(y)$.
\end{proof}

\begin{lemma}[Universal $B$-transfer diamond]\label{lem:b-diamond}
Let $x>0$, let $e=1-\alpha(x)$, put $a=J(x)$ and $g=1-e$, and factor
\[
  3a+1-2e=2^jJ(y),\qquad j\ge2,
\]
so $y=B(x)$. Define
\[
 P=Q_1^e(a),\quad U=Q_2^e(a),\quad
 D=Q_j^g(J(y)),\quad C=Q_{j-1}^g(J(y)),\quad F=A(U).
\]
Then $D=A(P)$ and $C=B(P)=B(U)$. If $X$ is the common child of $D,C$ and
$Y$ is the other child of $C$, then
\[
  B(F)=Y.
\]
Consequently an obligation in the $B$-selecting phase transfers a continuing
\DRAW\ to the adjacent frame $\{D,C\}$; when $j=2$ the transferred frame is
again an obligation in the opposite phase.
\end{lemma}

\begin{proof}
The formulas for $D$ and $C$ are the signed boundary formula and the
shared-child identity. The remaining equality is a direct substitution in the
three cases $j=2$, $j=3$, and $j\ge4$: the constant suffix of $D$ determines
its common child $X$, while the one-step longer factor state $F$ has the same
raw side $Y$. For the outcome assertion, a common child of a \DRAW\ cannot be
\LOSS. If the exponent-one member is \DRAW, the frame $\{D,C\}$ already
contains a \DRAW. In the parent form of an obligation, the only remaining
orientation has the exponent-two member \DRAW; if $C$ were \WIN, outcome
recursion through the displayed diamond would force $Y$ both \LOSS\ and a
child of the \DRAW\ state $F$, impossible. Hence $C$ is \DRAW. For $j=2$,
$D,C$ are exactly the exponent-two/one pair over $J(y)$.
\end{proof}


\section{The closed routing bridge}\label{sec:closed-routing}

The local arithmetic in Sections 2--8 is deliberately left unchanged from the
audit-hardened reduction.  We now close the two global obligations that were
left open there.  Two negative regressions are retained throughout:

\begin{enumerate}[label=(N\arabic*)]
\item a smaller canonical coefficient is not, in the presence of a factor
      $3^k$, a smaller coefficient source;
\item a $B$ return after several canonical $A$ steps need not lie below the
      numerical source present before that streak.
\end{enumerate}

The proof below never uses either implication.  Numerical values displayed
along a canonical streak are routing cursors until a stated source comparison
promotes one of them to a ranked anchor.

\subsection{Permanent multi-$A$ regression}

\begin{lemma}[Pre-streak anchor counterexample]\label{lem:anchor-stress}
Let $x_0=20$ and $x_{i+1}=A(x_i)$.  Then
\[
 20\longmapsto30\longmapsto45\longmapsto68,
 \qquad B(68)=25>20.
\]
The phase carried by the canonical lift switches to the $B$-selecting phase
at $68$, and the corresponding signed valuation is $3$.
\end{lemma}

\begin{proof}
The three $A$ values are immediate.  Moreover
\[
 \alpha(20),\alpha(30),\alpha(45),\alpha(68)=1,0,1,1,
\]
whereas the lifted phase flips after each canonical $A$ continuation.  At
$68$ the carried phase is therefore $0\ne\alpha(68)$, and
\[
 3J(68)+1=616=2^3\,77,
 \qquad 77=J(25).
\]
Thus the return is $25$, not a value below the pre-streak source $20$.
\end{proof}

\subsection{The exact raw selector}

For $s>0$ and $e\in\{0,1\}$ define the ordinary source selected by phase $e$
by
\[
 \chi_e(s)=
 \begin{cases}
   A(s),&e=\alpha(s),\\
   B(s),&e=1-\alpha(s).
 \end{cases}
\]

\begin{lemma}[Raw-selector identity]\label{lem:raw-selector}
For every $s>0$ and phase $e$,
\[
 \boxed{R\!\left(Q_1^e(J(s))\right)=\chi_e(s).}
\]
\end{lemma}

\begin{proof}
Put $z=Q_1^e(J(s))=2J(s)-e$ and $t=R(z)$.  Since
$z\bmod2=e$, Lemma~\ref{lem:alt-lift}, applied to $z$, gives for some
$m\ge1$
\[
 3z+1=Q_m^{1-e}(J(t))=2^mJ(t)-(1-e).
\]
After substituting $z=2J(s)-e$ and dividing by $2$,
\[
 3J(s)+1-2e=2^{m-1}J(t).
\]
The unique factor-free source in this signed transition is therefore $t$.
Lemma~\ref{lem:source-boundary} says that this source is $A(s)$ in the
$A$-selecting phase and $B(s)$ in the opposite phase.  Hence
$t=\chi_e(s)$.
\end{proof}

The next lemma is the missing provenance statement behind the old
$A/B_2$ reset.

\begin{lemma}[Second selector across an ordinary side diamond]
\label{lem:second-selector}
Let $q>0$ be ordinary, put $y=B(q)$ and $g=\alpha(q)$.  Then
\[
 \boxed{B(A(q))=\chi_g(y).}
\]
In particular, if the phase at $q$ is $B$-selecting and its signed valuation
is $2$, the transferred obligation at $y=B(q)$ selects on its next source
step the exact common side $B(A(q))$.
\end{lemma}

\begin{proof}
Let $k$ be the alternating-suffix length of $A(q)$.  Ordinary means
$k\in\{1,2\}$.  The proof of Lemma~\ref{lem:side} gives the refined table
\[
\begin{array}{c|c|c|c}
 k&g& B(A(q))=A(y) & g=\alpha(y)\\ \hline
 1&1-(y\bmod2)&y\bmod4\in\{0,3\}&y\bmod4\in\{0,3\}\\
 2&y\bmod2&y\bmod4\in\{1,2\}&y\bmod4\in\{1,2\}.
\end{array}
\]
The formulas for $g$ follow directly by substituting the one-bit and two-bit
alternating suffixes in $A(q)=2^k y+t_{k,y\bmod2}$ and using the two-bit
definition of $\alpha$.  Thus $B(A(q))=A(y)$ exactly when
$g=\alpha(y)$; otherwise it is $B(y)$.  Lemma~\ref{lem:raw-selector}
identifies this with $\chi_g(y)$.

For the last assertion note that the signed $B$-valuation equals $k+1$.
Valuation $2$ therefore has $k=1$, so $q$ is ordinary automatically; the
transferred phase is $g=1-(1-\alpha(q))=\alpha(q)$.
\end{proof}

\begin{corollary}[Exceptional rows never form a $B_2$ recycle]
\label{cor:exceptional-no-b2}
If $q\bmod16\in\{1,3,12,14\}$, then the $B$-selecting signed valuation at
$q$ is at least $4$.
\end{corollary}

\begin{proof}
For an exceptional $q$, the alternating suffix of $A(q)$ has length at least
three.  The signed $B$-valuation is one more than that suffix length, hence is
at least four.
\end{proof}

\subsection{Retained data, two one-shot seeds, and phase alignment}

The earlier versions used two different finite-token layers.  That distinction
is essential and is restored here, but without the invalid pre-streak numerical
comparison of Lemma~\ref{lem:anchor-stress}.  A marked configuration retains

\[
  (s,\eta,M,\zeta,N,a;\ \text{routing data}).
\]

Here $s$ is the outer coefficient-source anchor; $M$ is a finite multiset of
certified outer proof-token heights; $\eta\in\{1,0\}$ is the unique outer
seed bit; $N$ is a finite multiset of certified inner proof-token heights;
$\zeta\in\{1,0\}$ is the inner seed bit; and $a$ is the retained numerical
anchor supplied by the currently active typed entry.  When the entry is a
free obligation its factor-free coefficient is $J(a)$.  In an attached
factor/lift task, $a$ is the nested source anchor inherited from the exact
parent/child geometry.  Temporary ordinary sources, phases, valuations and
tail exponents are routing data and are not silently promoted to $s$ or $a$.

Put
\[
 \Phi(M)=\mathop{\#}_{h\in M}\omega^h,
 \qquad
 \Psi(N)=\mathop{\#}_{h\in N}\omega^h,
\]
and order retained data lexicographically by
\begin{equation}\label{eq:v6-rank}
 \boxed{\Pi=(s,\eta,\Phi(M),\zeta,\Psi(N),a).}
\end{equation}
Both bits are ordered by $1>0$.  A step classified as an outer token
decrease preserves the preceding outer source $s$; a step classified as an
inner token decrease preserves $(s,\eta,M,\zeta)$.  Thus no decrease of
$\Phi(M)$ or $\Psi(N)$ can pay for an increase in an earlier retained
component.  The seed bits are governed by the following non-reseeding rule.

\begin{lemma}[Seed admissibility and non-reseeding]\label{lem:seed-policy}
At a new outer-source fibre, while $\eta=1$, the inner seed $\zeta$ is
unavailable.  The first exact finite WIN/LOSS token produced by any typed
entry diamond---boundary, loss-sibling, base-entry, or the first free
obligation---may be installed in $M$ only together with the unique strict
transition $\eta:1\to0$.  After that transition, $M$ changes only by
certified proper-descendant replacement.

At fixed $(s,\eta=0,M)$, the unique transition $\zeta:1\to0$ may install the
first exact finite token of the active attached inner task in $N$.  After
$\zeta=0$, $N$ changes only by certified proper-descendant replacement.  If
both seed bits have been spent, no unrelated finite state may ever be added
to either multiset merely because its outcome is known.  This restriction is
about \emph{ranked tokens}.  A typed fixed-fibre subroutine may use a finite
WIN/LOSS state of unrelated height as temporary routing data, provided it is
not inserted into $M$ or $N$ and the subroutine is proved well-founded for
every possible initial value of that temporary witness.  This is the well-founded-fibre distinction used later for attached fixed-fibre subroutines.

A strict decrease of $M$ may reset $\zeta,N,a$; a strict decrease of $N$ may
reset $a$.  With all token data fixed, $a$ changes only by a separately
proved strict source comparison.  In particular, the two seed transitions
cannot occur in the opposite order and neither can be repeated in an
unchanged earlier fibre.
\end{lemma}

\begin{proof}
This is exactly the lexicographic discipline of \eqref{eq:v6-rank}.  Since
$\eta$ precedes $\Phi(M)$, the transition $\eta:1\to0$ is strict even when it
installs an arbitrary finite first multiset $M$.  While $\eta=1$, allowing a
$\zeta$ transition would leave the earlier unused seed available and would
permit two incomparable initializations, so it is excluded by definition of
the marked relation.  Once $\eta=0$, $\zeta$ precedes $\Psi(N)$ and gives the
same one-shot payment for the first inner token.  After a seed is zero, adding
an unrelated token would increase the corresponding ordinal multiset and is
therefore not a permitted transition.  Only a strict decrease in an earlier
component permits later components to be reset.
\end{proof}

\begin{lemma}[Phase-aligned common side]\label{lem:phase-aligned-side}
Let $O(x,e)$ be an $A$-selecting obligation, so $e=\alpha(x)$.  Put
\[
 y=A(x),\qquad g=1-e,
\]
and let $b$ be the common side of its exponent-one/two boundary.  In the
canonical branch, whose next obligation is $O(y,g)$,
\[
 \boxed{b=\chi_g(y).}
\]
Consequently, if $g$ is $A$-selecting at $y$, then $b=A(y)$; if $g$ is
$B$-selecting, then $b=B(y)$.
\end{lemma}

\begin{proof}
With the notation of the obligation put
\[
 P=Q_1^e(J(x)),\qquad V=A(P).
\]
The canonical identity is
\[
 V=Q_1^g(J(y)).
\]
The common side is $b=B(P)=R(A(P))=R(V)$.  Lemma~\ref{lem:raw-selector}
therefore gives $b=\chi_g(y)$ exactly.
\end{proof}

\begin{lemma}[Finite-source front end]\label{lem:finite-front}
Suppose an $A$-selecting obligation $O(x,\alpha(x))$ is reached while the
ordinary source $x$ itself is already a retained finite WIN/LOSS token.
Before any later high return or free reset, one of the following occurs:
\begin{enumerate}[label=(\alph*)]
\item that token is replaced by a certified proper proof-tree descendant;
\item if $x$ is exceptional and its canonical LOSS child is $B(x)$, the
      route becomes an exact adjacent factor-free lift over this strict
      descendant token.
\end{enumerate}
The statement includes both possible outcomes, WIN or DRAW, of the common
side; no fresh finite endpoint of unrelated height is inserted.
\end{lemma}

\begin{proof}
Put $y=A(x)$ and let $b$ be the common side of the obligation.  As in the
universal obligation fork, $b$ is never LOSS.

If $x$ is LOSS, then $y$ is WIN and $h(y)\le h(x)-1$.  Since one child of
the WIN state $y$ is the non-LOSS state $b$, its other child $\ell$ is the
unique LOSS child.  Hence
\[
 h(\ell)=h(y)-1\le h(x)-2.
\]
Thus $x\mapsto\ell$ is strict, independently of whether $b$ is WIN or DRAW.

Now let $x$ be WIN and choose a canonical LOSS child $\ell_0$ of $x$.
If $\ell_0=y=A(x)$, then every child of $y$ is WIN; in particular the common
side $b$ is WIN and
\[
 h(b)\le h(\ell_0)-1=h(x)-2.
\]
So $x\mapsto b$ is strict.

It remains that $\ell_0=B(x)$.  If $x$ is ordinary, the ordinary side
identity makes
\[
 p=B(A(x))
\]
a child of the LOSS state $\ell_0$.  Therefore $p$ is WIN and
\[
 h(p)\le h(\ell_0)-1=h(x)-2.
\]
This strict token is retained before the obligation chooses any later
canonical or factor continuation.

If $x$ is exceptional, the unconditional exceptional child-pair identity
(Lemma~\ref{lem:exceptional-side}) gives the two children of $y=A(x)$ as an
exact adjacent pair
\[
 \{A(y),B(y)\}
 =\{Q_{m+1}^{\delta}(J(\ell_0)),Q_m^{\delta}(J(\ell_0))\}
\]
for some $m\ge1$ and phase $\delta$.  The token replacement
$x\mapsto\ell_0$ is strict because $h(\ell_0)=h(x)-1$, and the later
arithmetic is attached to this same token.  These cases are exhaustive.
\end{proof}

\begin{lemma}[Aligned obligation catches or spends its token]
\label{lem:aligned-catch}
Suppose $O(x,e)$ carries a retained WIN token
\[
 w=\chi_e(x).
\]
Then the next outcome-compatible macro either reaches an obligation whose
ordinary source is exactly the retained token $w$, or strictly replaces
$w$ by a proper descendant before any reset.

If $e$ is $B$-selecting, every valuation-two return reaches source $w$
exactly; valuations at least three enter an attached factor frame over $w$.
If $e$ is $A$-selecting, then $w=A(x)$.  A canonical continuation reaches
source $w$; every noncanonical continuation replaces $w$ by its unique LOSS
child before entering the attached factor/side constructor.
\end{lemma}

\begin{proof}
If $e$ is $B$-selecting, the selected ordinary source is by definition
$w=B(x)$.  Lemma~\ref{lem:b-diamond} transfers the DRAW to the exact adjacent
frame over $w$.  At valuation two this frame is the next obligation itself,
with ordinary source $w$; at larger valuation it is an attached factor frame
over the same retained token.

Now let $e=\alpha(x)$.  Then $w=A(x)$.  Let $d$ be the common side of
$O(x,e)$.  Lemma~\ref{lem:phase-aligned-side}, applied one level later,
places $d$ among the two ordinary children of $w$.  The state $d$ is never
LOSS.  If the obligation takes its canonical branch, the next ordinary
source is $A(x)=w$ exactly.  Otherwise $d$ is either DRAW or the WIN common
side of the factor branch.  Since the retained state $w$ is WIN and one of
its children $d$ is non-LOSS, its other child $\ell$ is the unique LOSS
child, so
\[
 h(\ell)=h(w)-1.
\]
Replace $w$ by $\ell$ before entering the corresponding attached side/factor
constructor.  Thus no noncanonical row can reset at equal token rank.
\end{proof}

\subsection{Closure of the seeded $A/B_2$ reset (former P1)}

\begin{proposition}[Seeded $A/B_2$ reset provenance]\label{prop:P1-closed}
Every outcome-compatible canonical $A$ continuation and every valuation-two
$B$ return from a typed free/aligned obligation satisfies the P1 acceptance
conditions.  More precisely, before a control reset to a fresh $A$-obligation
is permitted, one of the following has happened:
\begin{enumerate}[label=(\roman*)]
\item a retained numerical source has strictly decreased;
\item an available seed has installed an exact finite token according to
      Lemma~\ref{lem:seed-policy};
\item an existing retained proof token has been replaced by a certified proper
      proof-tree descendant; or
\item the route has become an attached factor/lift task over such a retained
      token, in which case it is no longer eligible for a free reseed.
\end{enumerate}
This statement covers canonical $A$-streaks of every outcome-compatible
length and all exceptional residue classes.  It never compares a final
$B$-return with the numerical cursor present before the streak.
\end{proposition}

\begin{proof}
First consider a free obligation whose ordinary source $x$ is the current
retained inner anchor and which does not yet carry an aligned finite token.
Such a free obligation is admissible only while a seed is available under
Lemma~\ref{lem:seed-policy}; once both seeds have been spent, all obligations
arrive through the aligned/attached cases below.  A $B$-selecting source
returns to $B(x)<x$, hence makes a genuine retained inner-source decrease.
At an $A$-selecting source the universal obligation fork has a common side
$b$ which is not LOSS.  If $b$ is DRAW, the standard source estimate
\[
 \rho(b)\le\frac{9x+1}{24}<x
\]
is a genuine source comparison against the factor-free retained source $x$
and gives (i).  If $b$ is WIN, install this exact boundary token using the
earliest admissible seed: $\eta:1\to0$ if this is the first finite token in
the outer fibre, otherwise $\zeta:1\to0$ with $\eta=0$ for the first inner
token of the current task.  No such installation is allowed if both bits are
zero.

Suppose first that the obligation then takes its factor branch.  The token
$b$ remains marked, and the factor fork is by definition an attached entry
over this exact token.  This is alternative (iv), so P1 stops here: it makes
no claim about how the factor fork is normalized and installs no additional
inner witness.  That strictly later task belongs to the factor/return closure
of Section~\ref{sec:P2closed}; if it eventually needs its first inner token,
Lemma~\ref{lem:seed-policy} governs the one possible use of $\zeta$.

Suppose instead that the obligation continues canonically.  With
$y=A(x)$ and $g=1-\alpha(x)$, Lemma~\ref{lem:phase-aligned-side} gives
\[
 b=\chi_g(y).
\]
Thus the next obligation is aligned with the token just installed.
Lemma~\ref{lem:aligned-catch} says that within this next macro either its
ordinary source becomes exactly $b$, or $b$ is strictly replaced before the
route becomes attached.  In particular there is no tokenless multi-$A$
streak after the first seed.

Once an ordinary obligation source itself is a retained finite token,
Lemma~\ref{lem:finite-front} applies to every $A$-selecting continuation and
forces a strict token replacement (or an attached exceptional constructor)
before any later high return.  Hence an arbitrarily long canonical $A$-streak
is a finite concatenation of strict token events and attached routing; it is
not an equal-rank walk through the numerical cursors
$x,A(x),A^2(x),\ldots$.  This is exactly the point at which the old
pre-streak estimate failed.

It remains to inspect a valuation-two $B$ return from a retained finite
source $q$.  Corollary~\ref{cor:exceptional-no-b2} excludes all four
exceptional residue classes.  Put $y=B(q)$ and let the transferred phase be
$g=\alpha(q)$.  If $q$ is LOSS, then $y$ is a WIN child with
$h(y)\le h(q)-1$, so $q\mapsto y$ is strict.  If $q$ is WIN, choose a
canonical LOSS child.  When that child is $y$, the same strict replacement
holds.  In the remaining orientation the canonical LOSS child is $A(q)$.
The ordinary side identity gives
\[
 c=B(A(q))\in\operatorname{moves}(A(q))\cap\operatorname{moves}(y),
 \qquad h(c)\le h(q)-2.
\]
Lemma~\ref{lem:second-selector} says that the transferred phase at $y$
selects this exact $c$.  Therefore the reset is entered already carrying the
strict descendant $c$; it is either caught by the aligned/source-token case
or treated as an attached task.  No temporary cursor is promoted.

Finally, if a canonical streak reaches a $B$-selecting phase only after
several $A$ steps, the preceding argument is still exhaustive: immediately
after the first outer seed the next obligation is aligned; within at most one
canonical step its source catches the retained token unless that token has
already decreased.  Every subsequent $A$-obligation whose source is finite
spends a descendant by Lemma~\ref{lem:finite-front}.  Thus every
outcome-compatible streak length is covered without a numerical bound on the
pre-streak source.  The mandatory regression
$20\to30\to45\to68$, $B(68)=25>20$, is therefore harmless: the proof never
uses the comparison $25<20$ (which is false).
\end{proof}


\section{Semantic closure of factor, return and attached routing}
\label{sec:P2closed}

The remaining task is to prove that an attached token cannot disappear when
the arithmetic passes through a factor fork, a high return, a marked tail or
a short terminal lift.  We state the local facts in the form needed for the
global rank.  For audit provenance, the same arithmetic identities also occur with longer
case-by-case derivations in Sections 91--135 of the frozen proof ledger
\texttt{docs/verified-results.md} at repository revision
\nolinkurl{9d12789de3dbd841851f3faff1a237018d43879d}.  The proofs below reproduce
the incidences and rank payments actually used, so that ledger is a
cross-check rather than a logical black box.  Sections 137--138 of the ledger
are not used here; the audited entry argument below replaces them.

\begin{lemma}[Finite-source export]\label{lem:finite-export}
Let $O(x,\alpha(x))$ be an $A$-selecting obligation, let its common side be
\WIN, and suppose the ordinary source $x$ is finite.
\begin{enumerate}[label=(\alph*)]
\item If $x$ is \LOSS, the canonical branch exports an actual \LOSS token
      at least two proof levels below $x$ before any later high return.
\item If $x$ is \WIN and either $x$ is ordinary or $A(x)$ is the
      canonical \LOSS child, the return carries a \WIN token of height at
      most $h(x)-2$.
\item In the sole remaining orientation, $x$ is exceptional and $B(x)$ is
      its canonical \LOSS child; the returned adjacent pair is an exact
      factor-free lift over that token, whose height is $h(x)-1$.
\end{enumerate}
\end{lemma}

\begin{proof}
Put $y=A(x)$ and let $b$ be the common \WIN side of the obligation.  If
$x$ is \LOSS, then $y$ is \WIN with $h(y)\le h(x)-1$.  Since one child of
$y$ is the \WIN state $b$, its other child is \LOSS and has height at most
$h(x)-2$.

If $x$ is \WIN, choose its canonical \LOSS child $\ell_0$.  When $x$ is
ordinary, Lemma~\ref{lem:side} makes $p=B(A(x))$ a child of $\ell_0$ in the
orientation where $\ell_0=B(x)$; if $\ell_0=A(x)$, $p$ is directly a child
of that LOSS state.  In either case $p$ is \WIN and
$h(p)\le h(x)-2$.  When $x$ is exceptional and
$\ell_0=B(x)$, Lemma~\ref{lem:exceptional-side} gives the two returned
children as adjacent lifts over $\ell_0$.  These are exactly the three rows.
\end{proof}

\begin{lemma}[Generic lower factor forks spend the retained mark]
\label{lem:factor-fork-pay}
Consider the factor alternative of a typed $A$-selecting obligation, with the
finite provenance supplied by that alternative.  Then every lower factor
exponent pays or remains attached before control can restart.

For lower exponent $r=1$, let $b$ be the WIN factor source, let $V$ be the
retained LOSS state from the factor alternative and put $D=B(V)$.  If $y$ is
the ordinary child of $b$ selected by the factor phase, the exact twin identity
is
\[
 \{D,y\}=\{A(b),B(b)\},\qquad D\ne y.
\]
Hence $D$ is WIN, $y$ is the actual LOSS child of $b$, and the row either
exposes a boundary WIN at least two proof levels below $b$ or transfers the
DRAW to an adjacent frame over this same LOSS token.

For every $r\ge2$, let $x$ be the finite ordinary source from which the
factor fork was generated.  Then $x$ is ordinary and the factor common side is
$b=B(A(x))$, a child of both $A(x)$ and $B(x)$.  Splitting the actual outcome
of $x$ gives respectively: a strict canonical coefficient-source edge for the
same DRAW state; a common WIN child of a canonical LOSS child with height at
most $h(x)-2$; or the LOSS state $A^2(x)$ with height at most $h(x)-2$.
\end{lemma}

\begin{proof}
In the $r=1$ factor alternative the standard twin calculation gives
\[
 \{D,y\}=\{A(b),B(b)\},\qquad D\ne y.
\]
Since $V$ is LOSS, its
child $D=B(V)$ is WIN.  The factor source $b$ is also WIN, so its other child
$y$ is necessarily LOSS and $h(b)=h(y)+1$.  If the factor phase is
A-selecting, the common side of the exponent-one/two frame is a child of $y$
and therefore is WIN with height at most $h(b)-2$.  A continuing DRAW member
of that frame points to this lower WIN boundary.  If the phase is B-selecting,
Lemma~\ref{lem:b-diamond} transfers the DRAW to the exact adjacent frame over
the selected source $y$, retaining the actual LOSS token.  No unproved
finite outcome is introduced.

For $r\ge2$, the signed transition at $A(x)$ is B-selecting, so the selected
source is $b=B(A(x))$.  The four exceptional residue classes would make this
phase A-selecting with valuation one; hence this row is ordinary.  By
Lemma~\ref{lem:side}, $b$ is a child of both ordinary children of $x$.

If $x$ is DRAW, no new game state is inferred: the \emph{same actual DRAW
state} $x$, written in its canonical constant-tail coordinates, has coefficient
source $\rho(x)<x$ by Lemma~\ref{lem:source-bound}.  Replacing the retained
inner numerical anchor $x$ by its genuine canonical source is therefore a
strict source edge.

If $x$ is WIN, choose a canonical LOSS child of $x$.  The common state $b$ is
a child of it, hence is WIN and $h(b)\le h(x)-2$.  If $x$ is LOSS, both
ordinary children are WIN; the factor row already makes $b$ WIN, so the other
child $A^2(x)$ of $A(x)$ is LOSS and has height at most $h(x)-2$.  These are
all outcome orientations and every finite-token comparison is a genuine
parent/child comparison.
\end{proof}

\begin{lemma}[High-return token descent]\label{lem:high-return-pay}
Assume the high-return guard, including the exact obligation
$O(c,\alpha(c))$ and the finite return provenance supplied by the preceding
factor scan. A high return can never be a pure reset.
\begin{enumerate}[label=(\alph*)]
\item In each short return row $v=5,6$, the exact finite return token
      $c=Q_{v-3}^g(3J(t))$ is already marked with its proved finite outcome.
      The common-side split either replaces a retained token by a proper
      descendant or enters the exact exceptional lift over its canonical
      lower LOSS token.
\item For every $v\ge7$ the return guard supplies finite WIN states
\[
 u=Q_{v-1}^g(J(t)),\qquad c=Q_{v-3}^g(3J(t)),
\]
with $A(u)$ LOSS. Put $p=B(A(u))=A(c)$ and let $b=B(p)$. Then
\[
 h(p)\le h(u)-2,\qquad h(b)\le h(c)-2.
\]
Thus the pair replacement $(u,c)\mapsto(p,b)$ is strict before either
continuation of $O(c,\alpha(c))$ is followed.

If the obligation takes its canonical continuation, the next ordinary
source is exactly $p=A(c)$, so the strict token $p$ is inherited and no
marked-tail task is installed. If it takes its factor alternative, there are
further finite states $q,y$ such that
\[
 (p,b)\longmapsto(q,y)
\]
is componentwise strict, the continuing DRAW lies in the exact frame
\[
 \{Q_3^g(J(b)),Q_2^g(J(b))\},
\]
and, with $D=v-6$ and $C=27J(t)$,
\[
 b=Q_D^g(C),\qquad
\begin{array}{c|cc}
 v&q\text{ WIN}&y\text{ LOSS}\\ \hline
 7&B(b)&A(b)\\
 8&A(b)&B(b)\\
 v\ge9&A(b)&B(b).
\end{array}
\]
Hence only the factor alternative enters the marked-tail control, and it
does so after a strict pair-token replacement.
\end{enumerate}
\end{lemma}

\begin{proof}
For $v=6$ the common side is $b=B(A(c))$. If this side is DRAW, the
constant-tail state $c$ cannot be WIN: its canonical LOSS child would also
have $b$ as a child and would force $b$ WIN. Hence $c$ is LOSS. The WIN child
$A(c)$ has the DRAW child $b$, so its other child $A^2(c)$ is LOSS with
$h(A^2(c))\le h(c)-2$.

For $v=5$ the common side is $b=A^2(c)$. If $b$ is DRAW and $c$ is LOSS,
then $A(c)$ is WIN and its other child $B(A(c))$ is LOSS at least two levels
below $c$. If $c$ is WIN, then $A(c)$ is not LOSS because it has the DRAW
child $b$; therefore $B(c)$ is the canonical LOSS child of $c$. In the
ordinary row $B(A(c))$ is a WIN child of this LOSS token, while in the
exceptional row Lemma~\ref{lem:exceptional-side} gives the exact adjacent
lift over it. If the common side is WIN instead, Lemma~\ref{lem:finite-export}
applies because the statement explicitly retains the finite outcome of $c$.
Thus both short rows pay or remain attached to an already lower token.

Now let $v\ge7$. Direct constant-tail reduction gives
\[
 p=B(A(u))=A(c).
\]
The return guard makes $A(u)$ LOSS, so $p$ is its WIN child and
$h(p)\le h(u)-2$. The state $c$ is ordinary in this range. Since $p=A(c)$
is WIN, the other child $B(c)$ is LOSS, and
\[
 b=B(p)=B(A(c))
\]
is a child of that LOSS state. Hence $b$ is WIN and
$h(b)\le h(c)-2$. This proves the first strict pair replacement before the
obligation branch is selected.

In the canonical alternative of $O(c,\alpha(c))$, the next source is
$A(c)=p$ exactly. Thus the already lower token is inherited and the route
returns directly to an obligation over $p$.

In the factor alternative, $p$ is ordinary. Because $p$ and its child $b$
are WIN, the other child of $p$ is LOSS. Its ordinary common child
$q=B(A(p))$ is WIN and satisfies $h(q)\le h(p)-2$. The factor phase selects
the other child $y$ of the WIN state $b$; since $q$ is already its WIN child,
$y$ is the unique LOSS child and $h(y)=h(b)-1$. Hence
$(p,b)\mapsto(q,y)$ is strict in the multiset order.

Finally direct substitution gives $b=Q_{v-6}^g(27J(t))$. The signed
transition at $b$ selects $A(b)$ at $v=7$, $B(b)$ with valuation at least
three at $v=8$, and $B(b)$ with valuation two for every $v\ge9$. In the
factor alternative the selected state is exactly the LOSS token $y$ and the
other child is the WIN token $q$, yielding the displayed orientation table.
Thus the factor branch, and only that branch, enters the marked-tail frame
with an actual continuing DRAW and with the required tokens already marked.
\end{proof}

\begin{lemma}[Every long marked tail pays its LOSS token]
\label{lem:marked-tail-pay}
Let $C$ be positive and odd, let $D\ge3$, and put
\[
 b=Q_D^g(C),\qquad q=A(b),\qquad y=B(b).
\]
Suppose the marked high-return constructor retains the exact outcomes and
DRAW frame
\[
 b,q\text{ \WIN},\qquad y\text{ \LOSS},\qquad
 \{Q_3^g(J(b)),Q_2^g(J(b))\}\text{ contains a \DRAW}.
\]
Put
\[
 r=\begin{cases}
    B(q)=B(y),&D=3,\\
    B(q)=A(y),&D\ge4.
   \end{cases}
\]
Then $r$ is a \WIN child of the retained LOSS state $y$, so
\[
 h(r)\le h(y)-1,
\]
and the complete marked tail may replace $y$ by $r$ before any inner reset.
\end{lemma}

\begin{proof}
For $D=3$, the exponent-two/one common-child identity gives
$B(q)=B(y)$.  For $D\ge4$, Lemma~\ref{lem:long-tail} gives
$B(q)=A(y)$.  In either row $r$ is a child of the LOSS state $y$, hence is
WIN and has strictly smaller height.

The attachment survives the complete two-level scan by exact coordinates,
not by a comparison of canonical coefficients.  At $D=3$ the first raw
member is $Q_M^g(J(r))$ for its actual suffix exponent $M\ge2$.  For every
$D\ge4$ it is exactly
\[
 S=Q_1^g(J(r)).
\]
If the signed branch is taken, its common side is an exact lift over the
phase-selected ordinary child of $r$; if neither first exit is DRAW, the
second bridge is the obligation over this same $S$ and its transferred
source is again an exact lift over that selected child.  At the three short
endpoints $D=3,4,5$ these states coincide with the first/second signed and
raw exits of the original factor scan over $r$.  Hence no branch introduces
a source unrelated to $r$.

Therefore the strict replacement $y\mapsto r$ may be performed before the
raw/signed choice, after which all later routing data are allowed to reset:
the multiset rank has already decreased.  The displayed identities prove the exact attachment and use no source
comparison on a factorful cursor.  They agree with the marked-tail identities
of the frozen local ledger.
\end{proof}

\begin{lemma}[The two shortest marked tails]
\label{lem:short-tail}
Let $C$ be positive and odd, put $b=Q_D^g(C)$, and assume the exact marked
frame contains a continuing DRAW.

For $D=1$ the retained orientation is
\[
 q=B(b)\text{ WIN},\qquad y=A(b)\text{ LOSS}.
\]
If
\[
 9J(b)+1-2g=2^jJ(w),
\]
then $j\ge2$; moreover $w=A(y)$ when $j=2$ and $w=B(y)$ when
$j\ge3$. Hence the $D=1$ row always replaces the LOSS token $y$ by an
actual WIN child before any inner reset.

For $D=2$ the retained orientation is
\[
 q=A(b)\text{ WIN},\qquad y=B(b)\text{ LOSS}.
\]
If $y=0$, the returned arithmetic has coefficient source zero and is routed
by Lemma~\ref{lem:zero-source}.  If $y>0$, the first factor numerator has
valuation exactly one and its returned source has the exact form
\[
 w=Q_m^\delta(J(y)),\qquad \delta=1-g,\qquad m\ge2.
\]
Thus the row is not a free restart: it is a constant-tail lift over the
already retained LOSS token $y$.
\end{lemma}

\begin{proof}
For $D=1$, $b=2C-g$. Direct substitution gives
\[
9J(b)+1-2g=
\begin{cases}
2(27C+5),&g=0,\\
2(27C-5),&g=1.
\end{cases}
\]
Since $C$ is odd, the parenthesized integer is even, so $j\ge2$.  Moreover,
with $y=A(b)=3C-g$ the same identity is exactly
\[
 9J(b)+1-2g
 =2\bigl(3J(y)+1-2(1-g)\bigr).
\]
The expression in parentheses is the ordinary signed source transition at
$y$ in phase $1-g$.  If its valuation is one (total $j=2$) it selects
$A(y)$; if its valuation is at least two (total $j\ge3$) it selects $B(y)$.
Thus $w=A(y)$ for $j=2$ and $w=B(y)$ for $j\ge3$. Both are actual children
of the retained LOSS state $y$ and are therefore WIN proper descendants. In
particular the apparent $j=1$ lift row sometimes listed in a coarse
trichotomy is arithmetically empty here.

For $D=2$, $b=4C-g$ and
\[
9J(b)+1-2g=
\begin{cases}
2(54C+5),&g=0,\\
2(54C-5),&g=1.
\end{cases}
\]
The parenthesized integer is odd, so the valuation is exactly one.  Let
$z=3C-g$; by the exponent-two/one boundary identity $y=R(z)$.  The returned
source $w$ is determined by
\[
 J(w)=54C+5-10g,
\]
so $w=18C+1$ for $g=0$ and $w=18C-2$ for $g=1$.  Its canonical odd
coefficient is
\[
 \kappa(w)=\begin{cases}
 \odd(9C+1),&g=0,\\
 \odd(9C-1),&g=1.
 \end{cases}
\]
Lemma~\ref{lem:alt-lift}, applied to $3C$ in the first row and to $3C-1$ in
the second, gives $\kappa(w)=J(y)$, so no factor of three is hidden in the
coefficient of $w$.

When $y>0$, let $\lambda$ be the maximal alternating-suffix length of $z$.
For $g=0$ the alternating-lift identity is
$3z+1=2^\lambda J(y)$, whence
\[
 w=2(3z+1)-1=Q_{\lambda+1}^{1}(J(y)).
\]
For $g=1$ the same identity reads $3z+2=2^\lambda J(y)$, whence
\[
 w=2(3z+2)=Q_{\lambda+1}^{0}(J(y)).
\]
Thus $\delta=1-g$ and $m=\lambda+1\ge2$.  If $y=0$, the complete-word row of
Lemma~\ref{lem:alt-lift} gives a source-zero constant-tail state instead.
\end{proof}

\begin{lemma}[LOSS-anchored lift and factor module]
\label{lem:loss-anchored-lift}
Let $y$ be an already retained finite LOSS token.  Suppose a continuing DRAW
obligation or adjacent factor frame is attached, by exact coordinates, to a
constant-tail state
\[
 x=Q_m^d(3^kJ(y)),\qquad m\ge1,\quad k\ge0.
\]
Then the attached normalization has only two kinds of outcome-compatible
macro-exit:
\begin{enumerate}[label=(\alph*)]
\item $y$ is replaced by one of its actual WIN children, or by a certified
      descendant of such a child, and the proof-token rank is strict;
\item $y$ is retained and an exact tail/factor counter decreases.  No
      unrelated finite endpoint is inserted into a ranked token multiset.
\end{enumerate}
In particular the module never exits as a fresh free obligation at an
unrelated numerical cursor.
\end{lemma}

\begin{proof}
Both ordinary children
\[
 c_A=A(y),\qquad c_B=B(y)
\]
of the LOSS state $y$ are WIN and satisfy
$h(c_A),h(c_B)\le h(y)-1$.  While a constant tail has length at least three,
Lemma~\ref{lem:long-tail} removes one or two tail bits and merely multiplies
the odd coefficient by $3$; its underlying three-free source remains $y$.
The valuation-two transfer and the canonical $A$ continuation therefore stay
inside the same attached module while the exact tail counter decreases.  No
finite endpoint encountered during this countdown is inserted into a ranked
multiset merely because its outcome is known.

It remains to audit a factor boundary.  The universal current/next-level
coupling reduces its first factor-free signed source over the finite source
$y$ to
\[
 9J(y)+1-2f=2^vJ(t).
\]
The exact source trichotomy is
\[
 t=\begin{cases}
   Q_r^\epsilon(J(c_B))\quad(r\ge1),&v=1,\\
   A(c_A),&v=2,\\
   B(c_A),&v\ge3.
  \end{cases}
\]
For $v=1$ this is Lemma~\ref{lem:alt-lift} applied to $A(y)$; for
$v\ge2$ it is the ordinary source-boundary transition at $A(y)$.  Thus the
$v=1$ row is already an exact lift over the proper WIN child $c_B$, whereas
$v\ge2$ is an actual child of the proper WIN child $c_A$.  Before any new
normalizer is installed we may therefore make the strict retained-token
replacement $y\mapsto c_B$ in the first row or $y\mapsto c_A$ in the other
two rows.  The returned cursor itself is not inserted as a token.

The consecutive-factor identity is equally important: when the current
valuation is at least two, the next factor source is an exact lift over the
already displayed $t$; when it is one, the next factor source is exactly one
ordinary selected-source step from that first lift.  Hence a DRAW choosing a
raw rather than signed exit cannot bypass the strict child token.  If powers
of three remain at tail exponent at least two,
Lemma~\ref{lem:fixed-tail-factor} removes them one level at a time or exposes
one of its typed finite side diamonds.  At exponent one the same two-level
factor coupling is used instead; the reverse-factor lemma is never invoked
there.  Consequently equal-token routing has a decreasing exact natural
counter, and every genuine restart is preceded by the displayed child-token
replacement.  This is the LOSS-source provenance lemma of the older local
proof, rewritten here so that it is not hidden inside the global assembly.
\end{proof}

\begin{lemma}[The $D=2$ exact lift remains LOSS-anchored]
\label{lem:D2-pure}
The $D=2$ row of Lemma~\ref{lem:short-tail} is an instance of
Lemma~\ref{lem:loss-anchored-lift} with retained LOSS token $y$.  Therefore
it either makes a strict descendant-token transition before control restarts,
or stays in a finite exact tail/factor countdown over the same $y$.  It never
uses a new seed and never becomes a free obligation at the returned cursor
$w$.
\end{lemma}

\begin{proof}
For $y=0$ this is Lemma~\ref{lem:zero-source}.  For $y>0$ the preceding
calculation gives $w=Q_m^{1-g}(J(y))$ with $m\ge2$.  If the first factor
numerator is $N_1=2J(w)$, then consecutive factor levels satisfy
\[
 N_2=3N_1-2(1-2g)
    =2\bigl(3J(w)+1-2(1-g)\bigr).
\]
Thus the next signed source is exactly the ordinary source-boundary
transition of this lift.  Lemma~\ref{lem:loss-anchored-lift} applies from
that point and keeps $y$ as the retained provenance until a strict
child-token event occurs.  A later marked-tail return is therefore either
after that strict event or still inside the same decreasing attached
countdown; there is no unranked same-fibre reseed.
\end{proof}

\begin{lemma}[Post-payment shortest marked lift arithmetic]
\label{lem:attached-lift}
Consider the $D=3$ row of Lemma~\ref{lem:marked-tail-pay} after the strict
replacement $y\mapsto r$ has already been made. Thus
\[
 b=Q_3^g(C),\quad q=A(b)\text{ WIN},\quad y=B(b)\text{ LOSS},\quad
 r=B(q)=B(y)\text{ WIN}.
\]
Since $q$ and $r$ are WIN, its other child $\ell=A(q)$ is LOSS. The two
children of $\ell$ are WIN and are proper descendants of $q$.

The first raw exit of the marked factor scan is an exact lift
\[
 S=Q_m^g(J(r)),\qquad m\ge2.
\]
For an arbitrary A-selecting lift cursor $x=Q_m^g(a)$, with
$e=1-g=\alpha(x)$, the common side $d$ of $O(x,e)$ satisfies
\[
 d=\begin{cases}
 9a-g,&m=2,\\
 R(Q_1^g(9a)),&m=3,\\
 Q_{m-3}^g(9a),&m\ge4,
 \end{cases}
\]
and the source transition at $A(x)$ has respectively valuation $1$, at least
$3$, and $2$. Hence every same-provenance canonical recycle with $m\ge4$
has the exact update
\[
 (m,a)\longmapsto(m-3,9a),
\]
which preserves the underlying three-free source and lowers the natural lift
counter. Every noncanonical exit remains an exact raw/signed boundary over
$r$ or over an actual child of the transported lower pair. No temporary
finite endpoint is promoted to the retained token multiset.
\end{lemma}

\begin{proof}
The identity $S=Q_m^g(J(r))$, $m\ge2$, is the exact appended-suffix
calculation at the $D=3$ boundary. The three-row formula follows by applying
Lemma~\ref{lem:long-tail} to $A(x)$ and by using the shared exponent-two/one
boundary at $m=2,3$. For $m\ge4$ the tail has at least four constant bits,
so one canonical A/B$_2$ block removes three of them and multiplies the odd
coefficient by $9$, giving the displayed update.

The retained geometry is independent of the value of $m$. The WIN state $q$
has the WIN child $r$, hence $\ell=A(q)$ is LOSS; both children of $\ell$
are WIN proper descendants of $q$. Repeated common-child steps transport this
pair componentwise while the raw exponent is reduced by three. Therefore the
factor/signed exits stay on the exact boundary generated by $r$ and its actual
children. This lemma asserts provenance transport only.  The proof-height
payment at the terminal classes is proved separately after this arithmetic
transport has been established.
\end{proof}

\begin{lemma}[Attached $B$-selecting exponent-one lifts descend]
\label{lem:attached-b}
Let $r>0$, $d\in\{0,1\}$, and put
\[
 S=Q_1^d(J(r)),\qquad e=1-d.
\]
Then $e$ is the $B$-selecting phase at $S$.  If $k=B(S)>0$, its state source
satisfies
\[
 \boxed{\rho(k)<r}.
\]
Hence a $B$-selecting obligation attached to this exact exponent-one lift
transfers, by Lemma~\ref{lem:b-diamond}, either to the terminal source or to
an adjacent frame with a strictly smaller retained nested source.  It never
becomes a fresh free obligation at the numerical cursor $S$.
\end{lemma}

\begin{proof}
A one-bit tail of bit $d$ has A-selecting phase $d$, so $e=1-d$ is
B-selecting.  The universal contraction and source bounds give
\[
 k=B(S)\le\left\lfloor\frac{3S+1}{4}\right\rfloor,
 \qquad
 \rho(k)\le\frac{k-1}{6}\le\frac{S-1}{8}.
\]
Since $J(r)\le3r+2$, one has $S\le6r+4$, and therefore
$(S-1)/8<r$ for $r\ge2$.  For $r=1$ the two lifts are $S=10,9$;
their B-children are $7,3$, both of state source $0$.  Thus the inequality
holds in every positive case.  Lemma~\ref{lem:b-diamond} supplies the exact
transferred adjacent frame over $k$.
\end{proof}

\begin{lemma}[Every terminal macro pays a retained token]
\label{lem:terminal-pay}
Suppose a terminal row is reached by the post-payment $D=3$ lift transport of
Lemma~\ref{lem:attached-lift}.  Let $\ell=A(q)$ be the retained LOSS token
whose two WIN children form the transported lower pair; in the zero-recycle
endpoint the already retained marked LOSS token is $y$.  The common-child transport is componentwise: after every
three-bit recurrence both transported WIN states are proper descendants of
the corresponding incoming WIN states.  At the terminal barrier they are
the complete child set of one reconstructed LOSS state $Z_n$.

Every terminal exponent-one, exponent-two, or exponent-three restart is
therefore preceded by a strict replacement of retained provenance.  More
precisely:
\begin{enumerate}[label=(\alph*)]
\item the zero-recycle two-bit row $m=2$ either exposes a boundary WIN at
      least two proof levels below its boundary-WIN source, or enters an exact
      adjacent DRAW frame over that source's actual LOSS child one level
      lower;
\item the exponent-one row first reconstructs a canonical LOSS token $L$
      with $h(L)\le h(\ell)-2$ and only then enters its $B$-selecting gate;
\item every positive-recycle exponent-two or exponent-three row reconstructs
      a LOSS token $Z_n$ with $h(Z_n)<h(\ell)$ before the new inner task is
      installed;
\item in the sole zero-recycle exponent-three row, the existing marked LOSS
      token $y$ is replaced by its actual WIN child $r$ with
      $h(r)\le h(y)-1$.
\end{enumerate}
Consequently no terminal row is a pure control reset.
\end{lemma}

\begin{proof}
First isolate the zero-recycle two-bit row, which must not be hidden inside
the positive-recycle inverse-parent formulas.  Let
\[
 x=Q_2^g(a)\text{ WIN},\qquad e=1-g,\qquad O(x,e),
\]
and put
\[
 y=A(x),\quad c=B(x),\quad d=A(y),\quad p=B(y),
\]
where $d$ is the common side of the obligation.  If $d$ is DRAW, then $y$ is non-LOSS, so the WIN state $x$ has the
actual LOSS child $c$ and $h(c)=h(x)-1$.  If $x$ is ordinary, the ordinary
side identity makes $p=B(A(x))$ a child of the LOSS state $c$.  Hence $p$ is
WIN and $h(p)\le h(x)-2$.  The two children of $y$ are precisely $d$ and
$p$; since $d$ is DRAW and $p$ is WIN, outcome recursion now gives that $y$
is DRAW.  Thus $y\to p$ is the required lower boundary.  If $x$ is
exceptional, Lemma~\ref{lem:exceptional-side} gives the exact adjacent DRAW
frame over the actual LOSS child $c$.

If $d$ is WIN, apply the finite-source export lemma
(Lemma~\ref{lem:finite-front}) directly to the finite WIN source $x$ and the
A-selecting obligation $O(x,e)$.  It gives a certified proper descendant
before the continuation can restart, except in the unique exceptional
orientation in which the canonical LOSS child is $c=B(x)$; there the exact
exceptional child-pair identity retains an adjacent frame over $c$, with
$h(c)=h(x)-1$.  Thus the $m=2$, zero-recycle row always pays a genuine
descendant token before any restart.

The remaining componentwise claim is the exact inverse-parent calculation for
the shortest marked pair after at least one three-bit transport step (or for
the zero-recycle exponent-three row).  If the first raw lift exponent is $M=3h+n$ with
$n\in\{1,2,3\}$, one common-child recurrence replaces both coordinates by
actual children and lowers their proof heights componentwise; after $h$
recurrences the pair is one of the three terminal barriers.  In the
exponent-two barrier the two WIN states are exactly the children of
\[
 Z_2=Q_2^e(3^{2h-1}J(r))\quad(h\ge1),
\]
and in the positive exponent-three barrier they are the two children of
\[
 Z_3=Q_3^e(3^{2h-1}J(r))\quad(h\ge1).
\]
The LOSS height formula (one plus the maximum of its two child heights)
preserves the componentwise decrement.  In the exponent-one barrier, for
$h\ge2$ the two WIN states are exactly the children of
\[
 Z_1=Q_4^e(3^{2h-3}J(r)),
\]
while for $h=1$ this parent is the original $\ell$.  Hence
\[
 h(Z_1)\le h(\ell)-2h+2,
\]
and the analogous positive-recycle $Z_2,Z_3$ bounds are strict.  These are
exact state identities; no new finite outcome is inferred from a numerical
source label.

For the exponent-one row retain the terminal notation
\[
 a=3^{2h}J(r),\qquad h\ge1,\qquad X=Q_1^g(a),\qquad e=1-g.
\]
The reverse terminal barrier reconstructs a LOSS state $Z_1$.  Its upper
child is
\[
 U=Q_3^e(a/9)=A(Z_1)\quad\text{WIN}.
\]
The two children of $U$ are exactly
\[
 Q_2^e(a/3),\qquad Q_1^e(a/3).
\]
Choose a canonical LOSS child $L$ of this WIN state.  Then
\[
 h(L)=h(U)-1\le h(Z_1)-2\le h(\ell)-2,
\]
(the sharper bound in the positive-recycle rows is not needed here).  With
$c=a/3$ the new terminal source is $X=Q_1^g(3c)$ and its phase $e$ is
$B$-selecting.  Thus the old token has already been replaced strictly before
this gate is normalized.

For completeness, that gate has no hidden unbounded restart.  Put
\[
 W=B(Q_1^e(c))=B(Q_2^e(c)).
\]
Since $L$ is one of the displayed exponent-one/two states, $W$ is an actual
WIN child of $L$ and $h(W)\le h(L)-1$.  If
$\lambda=\operatorname{asl}(9c-g)$ is at least four, its transferred source
$R(9c-g)$ is strictly below $c$.  For $\lambda=1,2,3$ the exact short rows
are respectively: a factor-free lift over $R(W)<c$; a factor-free lift over
the lower WIN token $W$; or the ordinary child $A(W)$.  Hence all three
short rows occur only after the strict source/token payment already displayed.
They may therefore initialize later inner routing data, but none is a
same-rank reseed of a free obligation.

For positive-recycle exponent three, the reconstructed parent can be written
\[
 Z_3=Q_3^e(3^{2h-1}J(r))\quad(h\ge1),
\]
whose two terminal barrier children are WIN; the reverse construction gives
$h(Z_3)\le h(\ell)-2h<h(\ell)$.  For exponent two ($h\ge1$),
\[
 Z_2=Q_2^e(3^{2h-1}J(r))
\]
has the two reconstructed WIN barrier children and likewise
$h(Z_2)\le h(\ell)-2h<h(\ell)$.  These are strict replacements before the
new obligation/factor state is installed.

Finally, in the zero-recycle exponent-three row the marked configuration is
precisely the $D=3$ case of Lemma~\ref{lem:marked-tail-pay}; hence
$r=B(q)=B(y)$ is an actual WIN child of the already present LOSS token $y$.
Thus $h(r)\le h(y)-1$.  Every terminal row therefore pays a source/token
component before any reset.
\end{proof}

\begin{corollary}[The post-payment $D=3$ lift module is well-founded]
\label{cor:short-lift-closed}
After the strict marked-tail token edge at $D=3$, every outcome-compatible
shortest-row lift continuation either makes another strict retained
source/token transition or, with retained provenance fixed, decreases the
exact natural lift exponent until it reaches $m\in\{1,2,3\}$, where
Lemma~\ref{lem:terminal-pay} makes a strict retained-token transition before
any restart. Hence no equal-retained-data cycle is contained in this module.
\end{corollary}

\begin{proof}
Lemma~\ref{lem:attached-lift} gives
$(m,a)\mapsto(m-3,9a)$ for every same-provenance canonical row with
$m\ge4$ and transports the exact lower pair through all temporary factor
exits. Thus $m$ reaches $1,2$ or $3$ in finitely many recurrences unless a
strict source/token exit occurs earlier. At those barriers
Lemma~\ref{lem:terminal-pay} supplies the strict payment. The arithmetic
transport is proved before the terminal payment, so the dependency is
one-way rather than circular.
\end{proof}

\begin{proposition}[Semantic exhaustiveness of the marked router]
\label{prop:P2-closed}
Every actual outcome-compatible DRAW continuation entering the typed routing
system belongs to one of the semantic rows below.  In each row all finite
states inserted into $M$ or $N$ are states already present in the displayed
parent/child geometry; no unrelated finite endpoint is promoted merely because
its outcome is known.

\begin{enumerate}[label=\textnormal{(S\arabic*)},leftmargin=*]
\item \textbf{Boundary or loss-sibling countdown.}  A factor-free frame
\[
 \{Q_{m+1}^e(J(r)),Q_m^e(J(r))\},\qquad m\ge1,
\]
known to contain a DRAW uses Lemma~\ref{lem:long-tail} while $m\ge3$.
The exact tail exponent decreases.  At $m=1,2$ the row is respectively an
obligation/factor boundary or the typed three/two constructor of
Lemma~\ref{lem:hidden-parent}.  A loss-sibling entry of
Definition~\ref{def:entry-controls} carries, in addition, its exact finite
WIN/LOSS witness and the finite local counter $\tau$ supplied by the
constructing diamond.  The fixed-tail and exponent-one predecessor rows
remove one factor/tail level at each reverse normalization step; the
ordinary, exceptional and hidden-parent side rows are one-step exact
attachments.  Thus before leaving (S1) the witness is merely retained (or is
replaced by a proved descendant) while $\tau$ strictly decreases to a
boundary, free obligation, or attached factor/lift row.  No fresh token is
needed inside this countdown.

\item \textbf{Free/outer obligation.}  A free obligation before the first
finite token may use the outer seed only according to
Lemma~\ref{lem:seed-policy}; after $\eta=0$ a first inner witness may use
$\zeta$ once.  Every canonical continuation and valuation-two reset is
exactly one of Proposition~\ref{prop:P1-closed}'s source/seed/descendant/
attached exits.  A reverse edge to a new free obligation is forbidden unless
that proposition has already paid one of those components.

\item \textbf{Finite-source A obligation.}  If its ordinary source is an
outer or inner finite token, Lemma~\ref{lem:finite-front} exports a certified
proper descendant before a later high return.  In the exceptional orientation
the two children are the exact adjacent lifts over the canonical lower LOSS
token given by Lemma~\ref{lem:exceptional-side}.

\item \textbf{Lower factor fork.}  Its lower exponent is exactly $r=1$ or
$r\ge2$.  Lemma~\ref{lem:factor-fork-pay} gives, respectively, a lower
boundary/LOSS token or the source/WIN/LOSS trichotomy.  The only nonterminal
continuation is the typed high-return scan carrying the same finite
provenance.

\item \textbf{Short high return.}  The returned valuation is $v=5$ or
$v=6$, and the constructor carries the finite return token
$c=Q_{v-3}^g(3J(t))$ with its proved outcome.  Lemma~\ref{lem:high-return-pay}(a)
gives the exact common-side split: a certified proper descendant, or the
exceptional adjacent lift over the canonical lower LOSS token.  Hence neither
short row is a pure reset.

\item \textbf{Long high return.}  The returned valuation is $v\ge7$.  With
\[
 u=Q_{v-1}^g(J(t)),\qquad c=Q_{v-3}^g(3J(t)),
\]
the guard includes $O(c,\alpha(c))$ and the finite return provenance.
Lemma~\ref{lem:high-return-pay}(b) first makes the strict pair replacement
$(u,c)\mapsto(p,b)$.  The canonical alternative then returns directly to an
obligation at the already lower token $p$.  Only the factor alternative makes
the second strict replacement $(p,b)\mapsto(q,y)$ and enters the exact
marked-tail frame. Thus every long high return is strict, but only its factor
branch creates a marked tail.

\item \textbf{Marked tail.}  Its guard records a positive odd $C$, a phase $g$,
and an exact tail length $D$, with $b=Q_D^g(C)$ WIN and the exact marked
factor frame
\[
 \{Q_3^g(J(b)),Q_2^g(J(b))\}\text{ containing a \DRAW}.
\]
The retained WIN/LOSS pair is orientation-sensitive:
\[
\begin{array}{c|cc}
D&q\text{ WIN}&y\text{ LOSS}\\ \hline
1&B(b)&A(b)\\
2&A(b)&B(b)\\
D\ge3&A(b)&B(b).
\end{array}
\]
Lemma~\ref{lem:short-tail} proves that $D=1$ is itself strict: its first
factor source is an actual child of the retained LOSS token.  At $D=2$ the first source is an exact lift over the retained LOSS token
$y$. Lemmas~\ref{lem:loss-anchored-lift} and~\ref{lem:D2-pure} keep this
row attached: no unrelated finite endpoint is promoted, equal-token routing
only decreases its exact tail/factor counter, and every restart pays a strict
child-token edge.  For every $D\ge3$,
Lemma~\ref{lem:marked-tail-pay} replaces the retained LOSS token $y$ by its
actual WIN child
\[
 r=B(q)=B(y)\quad(D=3),\qquad
 r=B(q)=A(y)\quad(D\ge4),
\]
so $h(r)<h(y)$ before any new task is installed.

\item \textbf{Attached lift.}  There are two semantically distinct attached
lift origins, and they are not conflated.  The $D=2$ row is the exact lift over the already retained LOSS token $y$ of
Lemma~\ref{lem:short-tail}. It is normalized by
Lemma~\ref{lem:loss-anchored-lift}: while $y$ is fixed an exact tail/factor
counter decreases, and before any genuine restart $y$ is replaced by an
actual WIN child or certified descendant. It is not sent through the $D=3$
terminal-barrier proof.

The post-payment $D=3$ row has an exact raw lift $Q_m^g(J(r))$, $m\ge2$,
over the lower token $r$. Lemma~\ref{lem:attached-lift} gives its complete
three-row arithmetic.  The stable row $m\ge4$ performs only
$(m,a)\mapsto(m-3,9a)$, and the terminal rows are paid by
Lemma~\ref{lem:terminal-pay}.  Corollary~\ref{cor:short-lift-closed} closes
this module without any appeal to the $D=2$ row.

\item \textbf{Attached B lift.}  Every B-selecting exponent-one lift has the
form
\[
 S=Q_1^d(J(r)),\qquad e=1-d.
\]
Lemma~\ref{lem:attached-b} gives either terminal source or
$\rho(B(S))<r$.  Thus the selected source is a strict nested-source exit;
it is not a free reseed at the numerical cursor $S$.

\item \textbf{Post-payment $D=3$ terminal barriers.}  The terminal exponent is
exactly one, two, or three inside the shortest marked lift of (S8).
Lemma~\ref{lem:terminal-pay} treats all three.  The $D=2$ LOSS-anchored entry never uses this terminal lemma; it is already
covered by Lemmas~\ref{lem:loss-anchored-lift} and~\ref{lem:D2-pure}.
Exponent one
first reconstructs a lower LOSS token and then has only the symbolic suffix
rows $\lambda=1,2,3$ or $\lambda\ge4$; exponent two and positive-recycle
exponent three reconstruct a strictly lower LOSS token; the sole
zero-recycle exponent-three row replaces the existing LOSS token by its
actual WIN child.  Thus every terminal return pays before re-entering an
obligation.
\end{enumerate}

These rows cover every exponent and every returned tail length.  Therefore
P2 is closed: an upward control return is either preceded by a strict
retained source/token edge or is a finite same-retained-data route whose
remaining exponent is explicitly decreasing.
\end{proposition}

\begin{proof}
The coverage is by the exact parameters used to construct the controls, not
by a numerical cutoff.  Boundary and loss-sibling entries carry the true
finite local tail/factor counter $\tau$, and marked-tail controls carry their
true constant-tail exponent.  Factor forks split at $r=1$ versus $r\ge2$.
High returns split at the only two short valuations $5,6$ versus the uniform
range $v\ge7$.  Attached lifts split at $m=1,2,3,\ge4$.  Terminal rows split by
exponent $1,2,3$, and the exponent-one suffix is split by its actual
alternating-suffix length $\lambda=1,2,3,\ge4$.  These partitions are
disjoint except where two exact descriptions meet at a boundary, and in such
overlaps the certified conclusion is the same.

The parent/child incidences and rank payments are precisely the conclusions
of Lemmas~\ref{lem:factor-fork-pay}--\ref{lem:terminal-pay}.  A factor
alternative whose ordinary source is itself the retained anchor/token uses
the generic lower factor fork (S4).  A factor alternative at a temporary
short-lift cursor does \emph{not} use (S4): it remains inside (S8), keeps the
incoming $q/y$ provenance, and follows the exact finite lift/terminal
transport of Lemma~\ref{lem:attached-lift}.  Thus a temporary numerical
cursor is never promoted to a retained source merely to make a factor lemma
applicable.  The generic fork either pays a source/token component or reaches
(S5)--(S6).  A long return (S6) pays the pair-token edge before the marked
source is installed.  Its marked source then has an actual tail length $D$.
Rows $D\ge3$ pay the LOSS-token edge of (S7) before any raw/signed exit.
The $D=1$ row is itself strict. The $D=2$ row is not a free pure reset: it
remains LOSS-anchored by Lemma~\ref{lem:loss-anchored-lift}; equal-token
routing only decreases its exact tail/factor counter, and every restart first
pays a child-token edge.  Thus following the arrows gives the
closed cycle

\[
  \text{obligation}\to\text{factor fork}\to\text{high return}
  \to\text{marked tail}\to\text{attached lift}\to\text{obligation},
\]

but every traversal contains a strict retained edge before a new marked tail
is installed.  Inside the already post-payment $D=3$ lift, equal-retained-data
routing can additionally decrease the displayed natural lift exponent before
returning to the obligation.  The $D=2$ LOSS-anchored route is finite at fixed token/counter data and
cannot restart without the strict child-token event of
Lemma~\ref{lem:loss-anchored-lift}. Terminal exits in the $D=3$ module
are (S10) and pay a strict token first.  Therefore no finite raw cursor is ever promoted to
$M$ or $N$ merely because its outcome is known, and there is no semantic
``other'' row or unranked reverse edge.
\end{proof}

\section{Well-foundedness and the no-DRAW theorem}\label{sec:global}

\begin{lemma}[Well-founded-fibre principle]\label{lem:fibre-principle}
Let a transition relation carry a value in a well-order.  Suppose every
transition weakly decreases that value and the relation restricted to every
fixed value is well-founded.  Then the full relation is well-founded.
\end{lemma}

\begin{proof}
An infinite transition path would give a nonincreasing sequence in a
well-order.  It can contain only finitely many strict decreases, so from some
point onward the retained value would be constant.  The remaining infinite
path would lie inside one fixed fibre, contrary to the hypothesis.
\end{proof}

\begin{proposition}[Fixed-retained-data routing is well-founded]
\label{prop:fixed-fibre}
Fix the complete retained tuple
\[
 (s,\eta,M,\zeta,N,a).
\]
The routing relation that preserves every one of these values is
well-founded.
\end{proposition}

\begin{proof}
In a fixed fibre no seed transition, token replacement or source decrease is
available.  Proposition~\ref{prop:P2-closed} then leaves only finite routing
rows.

A boundary or loss-sibling tail has a finite exact exponent and decreases it
until an obligation/factor boundary.  A canonical A-side run cannot pass four
consecutive WIN side states (Corollary~\ref{cor:no-four-win}); before that it
reaches a B phase or one of the strict/attached exits.  A valuation-two B
reset cannot occur as a pure reverse edge: Proposition~\ref{prop:P1-closed}
requires a seed, source edge or token edge, all impossible in the fixed
fibre.  A larger B valuation enters a factor row.

A lower factor fork either leaves by a strict source/token edge or reaches a
high return.  Every high return is strict in a retained token by
Lemma~\ref{lem:high-return-pay}, so no high-return cycle survives in the
fixed fibre.  Consequently a marked tail cannot be \emph{re-entered} in a
fixed token fibre.  The $D=1$ row is itself strict, and every $D\ge3$ marked
row is strict by Lemma~\ref{lem:marked-tail-pay}.  The $D=2$ row may occur as an initial attached control after a high return,
but Lemma~\ref{lem:loss-anchored-lift} leaves no equal-fibre restart: while
its token is fixed the exact tail/factor counter decreases, and any return to
a fresh generic control is preceded by a strict replacement of that LOSS
token.

The post-payment $D=3$ shortest-lift module is the only remaining local module
with an unbounded displayed tail counter.  By Lemma~\ref{lem:attached-lift}
and Corollary~\ref{cor:short-lift-closed}, its sole same-provenance recurrence
is $m\mapsto m-3$ for $m\ge4$. Between two such updates the control moves
forward through finite raw-factor/terminal stages; at $m=1,2,3$ it meets a
terminal barrier and pays a token. Thus a fixed-fibre path is a finite
concatenation of decreasing natural counters and finite acyclic routing
segments. It cannot be infinite.
\end{proof}

\begin{proposition}[Well-founded marked normalizer]\label{prop:normalizer}
The complete marked routing relation is well-founded under the retained
projection $\Pi$ of \eqref{eq:v6-rank} together with
Proposition~\ref{prop:fixed-fibre}.
\end{proposition}

\begin{proof}
A strict outer-source transition lowers $s$ and may reset every later
component.  At fixed $s$, the unique transition $\eta:1\to0$ is strict and
may install $M$ and reset later data.  Once $\eta=0$, every change of $M$ is
a certified proper-descendant replacement and strictly lowers $\Phi(M)$ by
Lemma~\ref{lem:token-rank}; it may reset $\zeta,N,a$.

Only after $\eta=0$, at fixed preceding components, $\zeta:1\to0$ is a
one-time strict inner-token seed.  Thereafter every change of $N$ is a
certified descendant replacement
and strictly lowers $\Psi(N)$; it may reset $a$.  With both token layers
fixed, every change of the retained numerical anchor $a$ is a separately
proved strict source decrease.  Hence every transition weakly decreases
$\Pi$, and every transition that changes retained data decreases it
strictly.  The equal-$\Pi$ fibre is well-founded by
Proposition~\ref{prop:fixed-fibre}.  Lemma~\ref{lem:fibre-principle} now
proves the claim.
\end{proof}

\begin{proof}[Proof of Theorem~\ref{thm:main}]
Assume that a DRAW exists and choose the least coefficient source $s$ among
DRAW positions.  Start with the corresponding outer fibre and $\eta=1$.

If $s=0$, Lemma~\ref{lem:zero-source} performs a finite normalization to a
boundary or loss-sibling entry.  If $s>0$, follow a continuing DRAW through
long constant tails until exponent one or two.  Proposition~\ref{prop:boundary-reduction}
then gives a typed entry or the factor-free exponent-two base state; the
latter is handled by Proposition~\ref{prop:base-entry}.  The factorful
exponent-one case is handled by Proposition~\ref{prop:factorful-entry}; the
reverse-factor lemma is never used at exponent one.  At no point is a
decrease of a factorful canonical coefficient treated as a source decrease.

The first exact finite token may appear in a loss-sibling diamond, in the
factor-free base entry, or in the first free obligation.  Whichever typed
entry produces it, Lemma~\ref{lem:seed-policy} installs it in $M$ on the
unique transition $\eta:1\to0$; no earlier use of $\zeta$ is allowed.  If a
free obligation is reached before any finite endpoint, Proposition~\ref{prop:P1-closed}
either strictly decreases its retained source or produces exactly such a
seedable token.  Once $\eta=0$, the first additional finite witness of an
attached inner task may consume $\zeta$ once; after that, every free
$A/B_2$ reset is strict in an already declared source/token component or is
an attached route over existing provenance.  Every factor/high-return/marked-
tail/terminal route is exhaustive by Proposition~\ref{prop:P2-closed}.
Therefore all subsequent outcome-compatible normalization takes place in the
well-founded marked relation of Proposition~\ref{prop:normalizer}.

Every macro is justified by exact parent/child incidences and preserves an
actual DRAW continuation separately from its finite tokens.  A macro may use
a reverse parent as a proof transformation; it is not asserted to be a
forward play sequence.  Whenever it follows play, it uses an actual child.
A DRAW has no LOSS child and has at least one DRAW child, so whenever the
macro requires a continuing game branch a DRAW child exists.  Repeating the finite macros from the hypothetical DRAW
would therefore create an infinite path in the well-founded marked relation,
a contradiction.

Hence the conjugated game has no DRAW positions.  Proposition~\ref{prop:first-normal}
returns this conclusion to the original odd-integer game, with terminal
conjugated state $0$ corresponding to original state $1$.  Thus every
positive odd start has finite remoteness and optimal play reaches $1$.
\end{proof}


\section{Verification boundary and reproducibility}\label{sec:verification}

The theorem above is a human mathematical claim.  The accompanying Python
program checks the exact arithmetic identities used in Sections 2--9,
including the raw-selector and second-selector identities, both permanent
negative regressions, the exceptional no-$B_2$ assertion, and the declared
finite pure-control order.  It also checks the constant-tail, fixed-tail,
base-entry, zero-source and exceptional-side identities inherited unchanged
from version 5.0 over large finite ranges.

The finite program is supporting verification, not an infinite proof.  In
particular it does not assign outcomes to all integers and it does not replace
the proof-height arguments in Sections~\ref{sec:P2closed}--\ref{sec:global}.
The local high-return and terminal formulas used in the semantic closure are
cross-referenced, for audit provenance only, to Sections 91--135 of the frozen
repository proof ledger at revision
\nolinkurl{9d12789de3dbd841851f3faff1a237018d43879d}.  All statements needed for
the theorem are reproduced and proved in this manuscript.  The globally
audited but invalid entry statement formerly appearing as Section 137 of that
ledger is not a dependency of this manuscript.

No end-to-end Lean formalization is claimed.  The purpose of the machine files
is to make transcription and arithmetic regressions easy to attack
independently.

\appendix
\section{Permanent adversarial regressions}

The following examples are normative tests for any later revision.

\paragraph{Coefficient/source separation.}
For $s=1$, $J(1)=5$, $a=3J(1)=15$ and $e=1$, the common boundary child is
$R(44)=22$.  Its canonical coefficient is $11<15$, but $11=J(3)$, so its
coefficient source has grown from $1$ to $3$.  Therefore no factorful step may
use ``smaller coefficient'' as a synonym for ``smaller source''.

\paragraph{Pre-streak anchor separation.}
Lemma~\ref{lem:anchor-stress} gives
$20\to30\to45\to68$ and $B(68)=25>20$.  The proof of
Proposition~\ref{prop:P1-closed} deliberately compares only exact parent/child
proof tokens or the immediate typed inner source; it never compares $25$ with
$20$.

\section{Closed semantic inventory}\label{app:inventory}

\begin{center}
\small
\begin{tabular}{@{}p{.22\linewidth}p{.31\linewidth}p{.37\linewidth}@{}}
\toprule
Constructor & Exhaustive rows & Rank effect before restart \\
\midrule
loss-sibling entry & fixed-tail; exponent-one predecessor; side/hidden parent &
finite $\tau$ countdown while retaining the exact witness \\
free obligation & $A$ canonical; $B_2$; factor exit &
Prop.~\ref{prop:P1-closed}: source, seed, strict token, or attached entry \\
attached obligation & ordinary; exceptional exact lift &
retains existing token; factor/lift lemmas, never reseeds \\
factor fork & $r=1$; $r\ge2$ &
strict source/token or same-token high return \\
high return & $v=5$; $v=6$; $v\ge7$ &
strict descendant; for $v\ge7$ strict two-token multiset descent \\
marked tail & $D=1$; $D=2$; $D\ge3$ &
$D=1$ strict LOSS-child; $D=2$ LOSS-anchored countdown/child-token payment; $D\ge3$ strict $y\mapsto r$ \\
short lift & $D=2$ coupled entry; post-payment $D=3$ with $m\ge2$ &
$D=2$ stays LOSS-anchored and pays before restart; $D=3$ decreases $m$ by $3$ or pays a token \\
terminal & zero-recycle $m=2$; exponent $1$; positive $2/3$; zero $3$ &
strict descendant or exact lower-LOSS attached frame before restart \\
\bottomrule
\end{tabular}
\end{center}

\section{Frozen local-proof provenance map}

For reproducibility, the following map records where the same local incidence
calculations appeared in the research ledger.  It is an audit provenance map,
not an additional proof assumption: the theorem above uses the statements and
proofs reproduced in this manuscript.  At frozen revision
\nolinkurl{9d12789de3dbd841851f3faff1a237018d43879d}, the corresponding ranges
are:
\begin{itemize}[leftmargin=*]
\item Sections 91--95: obligation and lower factor-fork source/token exits;
\item Sections 97--103: short and long high-return provenance and the
      two-token descent;
\item Sections 104--120: transport of the lower token through raw/signed
      factor exits and exact lifts;
\item Sections 121--128: terminal reconstruction and the three finite gate
      lengths;
\item Sections 129--135: multiset token rank, fibre principle, terminal entry,
      the two short marked-tail rows, and the universal $D\ge3$ LOSS-token
      replacement.
\end{itemize}
Section 136 is useful as a historical assembly, but the present proof instead
uses Propositions~\ref{prop:P1-closed}, \ref{prop:P2-closed}, and
\ref{prop:normalizer}.  Sections 137--138 are explicitly excluded because the
former contained the audited coefficient/source shortcut.

\begin{thebibliography}{9}
\bibitem{Guy1986}
R. K. Guy, John Isbell's Game of Beanstalk and John Conway's Game of
Beans-Don't-Talk, \emph{Mathematics Magazine} 59(5) (1986), 259--269.

\bibitem{Guy1996}
R. K. Guy, Unsolved Problems in Combinatorial Games, Problem 42, in
\emph{Games of No Chance}, MSRI Publications 29, Cambridge University Press,
1996.

\bibitem{Althofer2024}
I. Alth\"ofer, M. Hartisch, T. Zipproth, Analysis of a Collatz Game and Other
Variants of the $3n+1$ Problem, \emph{Advances in Computer Games} (2024),
123--132. DOI: \url{https://doi.org/10.1007/978-3-031-54968-7_11}.

\bibitem{AlthoferPage}
I. Alth\"ofer, Collatz Problems: The Two-Player $3n+-1$ Game,
\url{https://althofer.de/collatz-prizes.html}.

\bibitem{Repo}
G. Pochuev, Optimal $3n\pm1$ Game: Proof and Verification Repository,
\url{https://github.com/Grisha-Pochuev/3n-plus-minus-1-game}.

\bibitem{Zenodo}
G. Pochuev, Optimal Play in the Two-Player $3n\pm1$ Game, Zenodo,
\url{https://doi.org/10.5281/zenodo.21844684}, 2026.
\end{thebibliography}

\end{document}
